% Options for packages loaded elsewhere
\PassOptionsToPackage{unicode}{hyperref}
\PassOptionsToPackage{hyphens}{url}
\PassOptionsToPackage{dvipsnames,svgnames,x11names}{xcolor}
%
\documentclass[
  10pt,
]{article}
\usepackage{amsmath,amssymb}
\usepackage{setspace}
\usepackage{iftex}
\ifPDFTeX
  \usepackage[T1]{fontenc}
  \usepackage[utf8]{inputenc}
  \usepackage{textcomp} % provide euro and other symbols
\else % if luatex or xetex
  \usepackage{unicode-math} % this also loads fontspec
  \defaultfontfeatures{Scale=MatchLowercase}
  \defaultfontfeatures[\rmfamily]{Ligatures=TeX,Scale=1}
\fi
\IfFileExists{lmodern.sty}{\usepackage{lmodern}}{}
\ifPDFTeX\else
  % xetex/luatex font selection
  \setmainfont[]{Latin Modern Roman}
  \setmonofont[Scale=0.85]{DejaVu Sans Mono}
\fi
% Use upquote if available, for straight quotes in verbatim environments
\IfFileExists{upquote.sty}{\usepackage{upquote}}{}
\IfFileExists{microtype.sty}{% use microtype if available
  \usepackage[]{microtype}
  \UseMicrotypeSet[protrusion]{basicmath} % disable protrusion for tt fonts
}{}
\makeatletter
\@ifundefined{KOMAClassName}{% if non-KOMA class
  \IfFileExists{parskip.sty}{%
    \usepackage{parskip}
  }{% else
    \setlength{\parindent}{0pt}
    \setlength{\parskip}{6pt plus 2pt minus 1pt}}
}{% if KOMA class
  \KOMAoptions{parskip=half}}
\makeatother
\usepackage{xcolor}
\usepackage[margin=2.4cm]{geometry}
\usepackage{longtable,booktabs,array}
\usepackage{calc} % for calculating minipage widths
% Correct order of tables after \paragraph or \subparagraph
\usepackage{etoolbox}
\makeatletter
\patchcmd\longtable{\par}{\if@noskipsec\mbox{}\fi\par}{}{}
\makeatother
% Allow footnotes in longtable head/foot
\IfFileExists{footnotehyper.sty}{\usepackage{footnotehyper}}{\usepackage{footnote}}
\makesavenoteenv{longtable}
\setlength{\emergencystretch}{3em} % prevent overfull lines
\providecommand{\tightlist}{%
  \setlength{\itemsep}{0pt}\setlength{\parskip}{0pt}}
\setcounter{secnumdepth}{-\maxdimen} % remove section numbering
% --- polished typesetting ---
\usepackage{microtype}                 % protrusion/expansion: better justification, fewer overfulls
\usepackage{fvextra}                   % line-breaking in code blocks
\fvset{breaklines=true,breakanywhere=true,fontsize=\small}
\usepackage{newunicodechar}
\renewcommand{\arraystretch}{1.18}     % a little row breathing room in tables
\setlength{\emergencystretch}{3em}     % last-resort to avoid overfull boxes
% explicit math font = Latin Modern Math (matches Latin Modern Roman body)
\ifdefined\setmathfont\setmathfont{Latin Modern Math}\fi
% text-mode Unicode the chapters use (so Latin Modern Roman renders them)
\newunicodechar{✓}{\ensuremath{\checkmark}}
\newunicodechar{✗}{\ensuremath{\times}}
\newunicodechar{→}{\ensuremath{\rightarrow}}
\newunicodechar{↔}{\ensuremath{\leftrightarrow}}
\newunicodechar{↑}{\ensuremath{\uparrow}}
\newunicodechar{↓}{\ensuremath{\downarrow}}
\newunicodechar{⇌}{\ensuremath{\rightleftharpoons}}
\newunicodechar{⇒}{\ensuremath{\Rightarrow}}
\newunicodechar{−}{\ensuremath{-}}
\newunicodechar{≈}{\ensuremath{\approx}}
\newunicodechar{⁺}{\textsuperscript{+}}
\newunicodechar{₂}{\textsubscript{2}}
\newunicodechar{₃}{\textsubscript{3}}
\newunicodechar{₄}{\textsubscript{4}}
\ifLuaTeX
  \usepackage{selnolig}  % disable illegal ligatures
\fi
\IfFileExists{bookmark.sty}{\usepackage{bookmark}}{\usepackage{hyperref}}
\IfFileExists{xurl.sty}{\usepackage{xurl}}{} % add URL line breaks if available
\urlstyle{same}
\hypersetup{
  pdftitle={VP Chemistry \& Electromagnetism},
  pdfauthor={Young jae Lee (ORCID 0009-0002-7535-8245), Independent researcher, Daegu, Korea},
  colorlinks=true,
  linkcolor={RoyalBlue},
  filecolor={Maroon},
  citecolor={Blue},
  urlcolor={RoyalBlue},
  pdfcreator={LaTeX via pandoc}}

\title{VP Chemistry \& Electromagnetism}
\usepackage{etoolbox}
\makeatletter
\providecommand{\subtitle}[1]{% add subtitle to \maketitle
  \apptocmd{\@title}{\par {\large #1 \par}}{}{}
}
\makeatother
\subtitle{Derived from a Single Anchor on the Jamming-Lattice Substrate
--- Foundations, Optics, and Applications}
\author{Young jae Lee (ORCID 0009-0002-7535-8245), Independent
researcher, Daegu, Korea}
\date{2026-06-14 · DOI (reserved) 10.5281/zenodo.20680541}

\begin{document}
\maketitle

{
\hypersetup{linkcolor=}
\setcounter{tocdepth}{2}
\tableofcontents
}
\setstretch{1.05}
\hypertarget{abstract}{%
\section{Abstract}\label{abstract}}

This volume consolidates the Volume-Particle (VP) account of
\textbf{electromagnetism, chemistry, and their engineering applications}
into a single reproducible whitepaper. Everything is derived from one
observational anchor --- the electron, treated as a rotational dent in
an infinitely rigid, fully packing ``jamming'' lattice --- using only
\(\pi\)-geometry and the electromagnetic inverse-square law. It is a
companion to the VP Theory (physics) whitepaper, from which it takes the
fundamental electromagnetic mechanism (§10, §14) as input.

Three results anchor the account. First, \textbf{light is not assumed
but emerges}: at the isostatic point of the jammed lattice the relaxed
shear modulus vanishes while the bulk modulus stays finite, so a single
longitudinal wave survives at \(c^2 = B/\rho\); a deterministic lattice
simulation reproduces this to 0.06\%. Second, on that substrate charge
is synchronized rotation, the Coulomb \(1/r^2\) follows from
three-dimensional flux conservation, the \(E/B\) geometry and
polarization are rotation, and the same wave spans conduction and
radiation by a single angle \(\sin\chi=\lambda/(mD)\); \textbf{and the
dynamic field equations, the Poynting energy flow, and a strict
light-cone causality are supplied in longitudinal form} (Chapter EM,
§EM.12, from the AQD axiomatic field dynamics). Third, chemistry follows
by sphere geometry and rotation --- bonding angles \(\arccos(-1/3)\),
the crystal-field ratio \(4/9\), the close-packing fraction \(0.7405\),
the nuclear magic skeleton, thermodynamics, equilibrium, kinetics,
solution and electrochemistry --- and the engineering applications
(catalysis, fertilizer, water, energy storage, new materials, and
\textbf{CO\(_2\) reduction}) are re-grounded on the \textbf{verified
\(d\)-band energy descriptor}, not on the refuted electron-amplitude
variable.

Every forced result has zero free parameters and is therefore
falsifiable; measured constants are labelled as calibration inputs and
never tuned per result; open problems --- most fundamentally the
electromagnetic coupling \(\alpha_{\mathrm{em}}\) itself --- are
admitted rather than disguised. All numbers are reproduced by
deterministic, standard-library Python modules (a 38-module chemistry/EM
core with a 176-row case ledger, plus a 9-module / 46-row applications
package) with verification harnesses (core 38/38 pass; applications 9/9
pass) and numeric-drift gates that check every displayed number against
canonical recomputation and module output (both pass, zero drift).

\textbf{Keywords:} Volume-Particle theory, jamming lattice, light
emergence, \(c^2=B/\rho\), electromagnetism, Poynting, causality,
Coulomb \(1/r^2\), propagation angle, VSEPR, crystal-field \(4/9\),
\(d\)-band catalysis, CO\(_2\) reduction, reproducibility,
falsifiability.

\hypertarget{introduction}{%
\section{Introduction}\label{introduction}}

\hypertarget{the-program-and-this-consolidation}{%
\subsection{The program and this
consolidation}\label{the-program-and-this-consolidation}}

The VP framework posits a vacuum built from infinitely rigid, fully
packing particles --- a jammed lattice. From this single substrate the
physics whitepaper derives the speed of light, the proton radius, and a
ladder of dimensionless constants. This whitepaper extends the program
into \textbf{electromagnetism, chemistry, and applications}, asking how
much of each follows with \emph{no adjustable parameters} once light has
emerged. This consolidated edition gathers what were previously separate
chapter files (EM, CH, CM, CC, CT) and the applications package (CA)
into one volume, and incorporates two completions made after the first
release: the \textbf{dynamic-field / Poynting / causality} addendum
(§EM.12) and the \textbf{CO\(_2\) electrocatalytic reduction} module
(§CA, re-grounded on the \(d\)-band descriptor).

\hypertarget{why-light-first}{%
\subsection{Why light first}\label{why-light-first}}

The organizing principle is sequencing: \textbf{light must emerge before
anything else follows.} Once the lattice supports a single longitudinal
wave at \(c^2 = B/\rho\) with linear dispersion \(\omega = cq\), a
wavelength exists; once a wavelength exists, the propagation angle, the
quantum diameter, refraction, and color follow; and once charge is
identified as synchronized rotation, the Coulomb law and the entire
chemical edifice follow. Chapter EM begins with the emergence of light
and develops electromagnetism (now including its dynamic and causal
structure); the chemistry chapters build on that foundation; the
applications chapters close the arc on verified physics.

\hypertarget{how-claims-are-graded-the-no-tuning-discipline}{%
\subsection{How claims are graded (the No-Tuning
discipline)}\label{how-claims-are-graded-the-no-tuning-discipline}}

\begin{itemize}
\tightlist
\item
  \textbf{{[}F{]} forced} results are dimensionless geometry or integers
  with nothing to adjust (the tetrahedral angle \(\arccos(-1/3)\), the
  ratio \(4/9\), the packing fraction \(\pi/(3\sqrt2)\)). They cannot be
  tuned, so they are \emph{refuted} if measurement disagrees.
\item
  \textbf{{[}CAL{]} calibration inputs} are measured constants (the
  fine-structure coupling, spectroscopic constants, standard potentials,
  \(d\)-band centers) fixed once and used consistently; quantities
  \emph{derived} from them and then matching measurement are
  predictions, not fits.
\item
  \textbf{{[}V{]} simulation-measured}, \textbf{{[}H{]} hypothesis}, and
  \textbf{{[}O{]} open} label what a deterministic simulation shows,
  what remains a working hypothesis, and what is honestly unsolved.
\end{itemize}

A central methodological result of this program is \textbf{negative}:
the early ``electron amplitude (fm)'' mechanism for catalysis,
bond-angle deviations, magnetic desalination, and black-copper power
generation was tested and \textbf{refuted}, and is replaced --- where a
working mechanism exists --- by the correct variable (the \(d\)-band
\emph{energy}, equilibrium thermodynamics, real magnetostatics,
Carnot-bounded storage). Chapter CT records this refutation in full; the
applications chapter is built on the replacement. Letting the refuted
mechanism go is what makes the surviving results trustworthy to build
on.

\hypertarget{reproducibility-as-a-first-class-requirement}{%
\subsection{Reproducibility as a first-class
requirement}\label{reproducibility-as-a-first-class-requirement}}

Every numeric claim is produced by a deterministic module using only the
Python standard library; determinism is verified by re-running and
comparing output hashes. Verification harnesses check execution,
determinism, dependency, and case-ledger integrity (chemistry/EM core:
38/38; applications: 9/9), and numeric-drift gates check that every
displayed number matches both the canonical closed-form recomputation
and the module output (both pass, zero drift). Nothing here is asserted
that the reader cannot re-run.

\hypertarget{boundary-with-the-physics-whitepaper}{%
\subsection{Boundary with the physics
whitepaper}\label{boundary-with-the-physics-whitepaper}}

The \emph{fundamental} electromagnetic mechanism --- light as a
rotational transverse wave, the single surviving longitudinal speed, and
the \(1/r^2\) Green function --- is established in the physics
whitepaper (§10, §14). This document consumes those results and develops
the chemical, optical, and dynamical consequences. The electromagnetic
coupling strength \(\alpha_{\mathrm{em}}\) is \textbf{not derived} in
this framework (physics §14.5); it is an honest open item, carried as
{[}CAL{]}.

\newpage

\hypertarget{electromagnetism-in-the-jamming-vacuum-from-the-quantum-lattice-to-light}{%
\section{Electromagnetism in the Jamming Vacuum: From the Quantum
Lattice to
Light}\label{electromagnetism-in-the-jamming-vacuum-from-the-quantum-lattice-to-light}}

\textbf{VP Chemistry \& Electromagnetism Whitepaper --- Chapter EM} DOI
(reserved): 10.5281/zenodo.20680541 · Companion to VP Theory (physics)
§10, §14, §15 Status: each result is graded \textbf{{[}F{]}} forced
(zero free parameters) · \textbf{{[}CAL{]}} calibration input (measured
constant, used consistently) · \textbf{{[}V{]}} simulation-measured ·
\textbf{{[}H{]}} hypothesis · \textbf{{[}O{]}} open. Every numeric claim
is reproduced by a deterministic, standard-library module (2× sha256
identical); module names are given inline.

\begin{quote}
One medium, one motion seen two ways. Charge is the bookkeeping of a
synchronized rotation; light is the transverse swing of that rotation
propagating on the single surviving longitudinal wave. Electricity, a
radio wave, and a γ-ray differ only by an angle.
\end{quote}

\hypertarget{em.0-governance-inherited-from-vp-theory-1}{%
\subsection{EM.0 Governance (inherited from VP Theory
§1)}\label{em.0-governance-inherited-from-vp-theory-1}}

This chapter inherits the No-Tuning / LOCK / Gate discipline of the
physics whitepaper. \textbf{The fundamental electromagnetic mechanism is
supplied by physics §10 (light, clock-free \(c\)) and §14 (force, the
\(1/r^2\) Green function). This chapter does not re-derive it; it takes
those results as input and develops the chemical and optical
consequences.} Where a quantity is a measured constant (notably the
fine-structure coupling \(\alpha_{\mathrm{em}}\)), it is marked
\textbf{{[}CAL{]}} and is never tuned to fit a downstream result.

\hypertarget{em.1-genesis-charge-as-synchronized-rotation-hv}{%
\subsection{\texorpdfstring{EM.1 Genesis: charge as synchronized
rotation
\textbf{{[}H{]}→{[}V{]}}}{EM.1 Genesis: charge as synchronized rotation {[}H{]}→{[}V{]}}}\label{em.1-genesis-charge-as-synchronized-rotation-hv}}

A quantum that receives energy can only store it as rotation, so energy
input forces rotation. This is grounded in established physics:
photon--photon collisions create matter (\(\gamma\gamma \to e^+e^-\);
\(E=mc^2\)), and photons carry helicity \(\pm 1\) --- circular
polarization is a \emph{rotating field}. \textbf{Electric charge is the
bookkeeping of one synchronized rotation} of the quanta; the sign
\(\pm\) is the handedness (CW/CCW) of the twist of a charge's rotation
relative to the surrounding common axis. The elementary charge is one
boundary-crossing rotation unit (universal), so the EM coupling carries
no nucleon-structure factor (contrast the mass/gravity sector, where
\(m_p/m_e = 6\pi^5\) and the \(82{+}7\) core appear).

\hypertarget{em.2-light-emergence-the-keystone-fv}{%
\subsection{\texorpdfstring{EM.2 Light emergence --- the keystone
\textbf{{[}F{]}/{[}V{]}}}{EM.2 Light emergence --- the keystone {[}F{]}/{[}V{]}}}\label{em.2-light-emergence-the-keystone-fv}}

Light is not assumed; it \emph{emerges} from the jammed lattice. The
chain (physics §SP):

\begin{enumerate}
\def\labelenumi{\arabic{enumi}.}
\tightlist
\item
  \textbf{{[}F, axiom{]}} Infinitely rigid, fully packing quanta ⇒ the
  vacuum is a jammed lattice.
\item
  \textbf{{[}H{]}→{[}V{]}} Rotation (temperature) drives the contact
  number \(z\) to the isostatic threshold \(z = 2d = 6\) (self-organized
  criticality).
\item
  \textbf{{[}V{]}} At the isostatic point the shear reserve is zero, so
  the relaxed shear modulus vanishes while the bulk modulus stays
  finite:
  \[ G_{\text{relaxed}} \to 0,\qquad B = \text{finite}\;\Rightarrow\; c^2 = \frac{B}{\rho}=K. \]
  The transverse acoustic wave dies; \textbf{one longitudinal speed
  survives.}
\item
  \textbf{{[}F{]}} A single speed gives linear dispersion and hence a
  wavelength: \[ \omega = c\,q \;\Rightarrow\; \lambda = c/\nu. \]
\end{enumerate}

\textbf{Direct demonstration (module \texttt{vp\_light\_emergence.py}).}
A 1-D chain of quanta, \(m\,\ddot u_n = k(u_{n+1}-2u_n+u_{n-1})\), is
given a right-moving Gaussian packet and integrated by leapfrog (no
randomness). The energy centroid propagates at
\(c_{\text{measured}} = 0.99940\) versus the closed form
\(c_{\text{theory}} = a\sqrt{k/m} = 1\) --- a \textbf{0.06 \%}
agreement. The exact dispersion \(\omega(q)=2\sqrt{k/m}\,|\sin(qa/2)|\)
tends to \(\omega = cq\) as \(q\to 0\) (\(v_{ph}/c \to 0.99999\)), and
the identity \(c^2=B/\rho\) is verified to machine precision. A 2-D
triangular lattice confirms that as the shear reserve \(\to 0\) the
transverse speed \(c_T \to 0\) while the longitudinal \(c_L\) survives
--- \textbf{light is the unique surviving longitudinal wave of the
jammed lattice.}

\textbf{Quantum diameter {[}F{]}.} One \(2\pi\) phase winding of the
carrier sets the rotational circulation length
\[ D = 2\pi\,a_{\text{phys}} = \frac{2\pi\lambda}{A} = 2\lambda_{C,e} = 6\pi^6 r_p, \qquad
   r_p = \tfrac{2}{\pi}\lambda_{C,p}, \] where \(r_p\) is the globally
stable fixed point of the stiffness-shell balance
\(\alpha x^{-5}=x^{-4}\Rightarrow x^*=2/\pi\) (physics §SP S4).
Numerically \(D = 4.8526\) pm. The module verifies \(D=6\pi^6 r_p\) to
within \textbf{+18.8 ppm}, which is exactly the headline residual
\(m_p/m_e = 6\pi^5\;(-19\text{ ppm})\): the small mismatch is the
\emph{known residual of the theory}, not an error.

\hypertarget{em.3-the-coulomb-sector-why-1r2-why-long-range-how-strong}{%
\subsection{\texorpdfstring{EM.3 The Coulomb sector: why \(1/r^2\), why
long-range, how
strong}{EM.3 The Coulomb sector: why 1/r\^{}2, why long-range, how strong}}\label{em.3-the-coulomb-sector-why-1r2-why-long-range-how-strong}}

\textbf{Why a force exists {[}F{]}.} Two sources' fields superpose; the
stored strain energy depends on separation and the medium relaxes it,
giving \(E_{\text{int}} = k\,q_1 q_2 / r\). Like sources repel, opposite
attract (a positive-definite strain structure).

\textbf{Why \(1/r^2\) {[}F{]}.} A fixed source flux is diluted over the
\(4\pi r^2\) shell (Gauss = Poisson Green's function). In \(d\)
dimensions the field falls as \(r^{-(d-1)}\), so \(d=3\) gives
\(1/r^2\). Module \texttt{vp\_electromagnetism.py} measures the exponent
directly: \(d=1,2,3,4 \to 0,-1,-2,-3\), confirming \(E\propto 1/r^2\) in
three dimensions as the geometric consequence of flux conservation.

\textbf{Why long-range / unscreened {[}H{]}.} Synchronization
spontaneously breaks a \emph{global} \(U(1)\) (the common phase), giving
a massless Goldstone mode and hence infinite range. (Were the \(U(1)\)
gauged, its breaking would make the mediator massive and short-ranged;
long-range EM requires the broken symmetry to be global --- a stated
departure.)

\textbf{Coupling magnitude {[}CAL{]}.} The elementary charge is one
universal rotation unit, so
\[ K_C = k_e e^2 = \alpha_{\mathrm{em}}\,\hbar c, \qquad
   k_e = \frac{\alpha_{\mathrm{em}}\hbar c}{e^2}. \] The module
reproduces the Coulomb constant \(k_e = 8.987552\times 10^9\) to
\(3\times 10^{-10}\)\%. \textbf{\(\alpha_{\mathrm{em}}\) is taken from
measurement; it is not derived here} --- an honest {[}CAL{]} item.

\hypertarget{em.4-the-eb-geometry-longitudinal-e-transverse-b-f}{%
\subsection{\texorpdfstring{EM.4 The E/B geometry: longitudinal \(E\),
transverse \(B\)
\textbf{{[}F{]}}}{EM.4 The E/B geometry: longitudinal E, transverse B {[}F{]}}}\label{em.4-the-eb-geometry-longitudinal-e-transverse-b-f}}

The same rotating medium supplies the vector sector. The organizing
statement: \textbf{the electric field is the force directed along the
light (propagation) direction; the magnetic field is the force directed
at \(90^\circ\) to it.} Both originate from one fact --- a charge's
rotational output is \emph{twisted} relative to the surrounding
synchronized quanta because its axis differs. A charge with axis aligned
to the common axis emits a pure radial (Coulomb) field and no magnetic
field. A moving charge tilts its axis by an amount \(\propto v/c\); the
propagation of that twisted region, transverse to the light direction,
is the magnetic field: \[ |\mathbf B| = \frac{v}{c}\,|\mathbf E|. \]
Module \texttt{vp\_electromagnetism.py} confirms
\(|\mathbf B|/|\mathbf E| = v/c\) exactly. Because \(\mathbf B\)
\emph{is} a rotation (a pseudovector), it is automatically
divergence-free, \(\nabla\!\cdot\mathbf B = 0\): the poles are the two
ends of one rotation axis and cannot be separated --- there are
\textbf{no magnetic monopoles}, and cutting a magnet yields two fresh
N/S pairs. The right-hand rule is the handedness of the rotation.

\hypertarget{em.5-polarization-as-rotation-helicity-f}{%
\subsection{\texorpdfstring{EM.5 Polarization as rotation (helicity)
\textbf{{[}F{]}}}{EM.5 Polarization as rotation (helicity) {[}F{]}}}\label{em.5-polarization-as-rotation-helicity-f}}

Circular polarization \emph{is} the rotating field; the two senses
(CW/CCW) are helicity \(\pm 1\). In Jones form the circular states are
\((1,\pm i)/\sqrt 2\). The module computes their Stokes vectors and
finds \(|S_3| = S_0\) with \(S_1 = S_2 = 0\) --- i.e.~pure rotation, no
linear component. Linear polarization is the equal superposition of the
two rotation senses. (The sign convention of \(S_3\) depends on the time
convention \(e^{\mp i\omega t}\); the invariant content is
\(|S_3| = S_0\) for circular light.)

\hypertarget{em.6-the-propagation-angle-conduction-radiation-fvp-prediction}{%
\subsection{\texorpdfstring{EM.6 The propagation angle: conduction ↔
radiation
\textbf{{[}F{]}/{[}VP-prediction{]}}}{EM.6 The propagation angle: conduction ↔ radiation {[}F{]}/{[}VP-prediction{]}}}\label{em.6-the-propagation-angle-conduction-radiation-fvp-prediction}}

Light is unambiguously an electromagnetic wave --- \emph{one} object.
Its longitudinal part (\(\nabla\!\cdot\mathbf u\), compression =
deficit) is the static Coulomb field; its transverse, propagating part
is light. What distinguishes electricity, a radio wave, and a γ-ray is
only the propagation angle \(\chi\), set by the wavelength through
physics §10.9:
\[ \sin\chi = \frac{\lambda}{mD},\qquad m = \lceil \lambda/D\rceil. \]
Short-\(\lambda\) γ-rays are near-longitudinal (\(\chi\to 0\),
penetrating); visible and radio waves are near-transverse
(\(\chi\to 90^\circ\)). Crucially the visible band is \textbf{near ---
not exactly --- transverse}: \(\chi \approx 89.9^\circ\), a small
\emph{falsifiable} departure from the textbook exactly-transverse wave.
Electric conduction is the extreme-longitudinal limit (\(\chi\to 0\)):
charges drift through a wire's internal quanta, each producing a
transverse twist that wraps the wire azimuthally as the magnetic field.
\textbf{Conduction and radiation are the same electromagnetic
phenomenon, differing only in angle (longitudinal vs transverse) and
medium (atomic-internal quanta vs free space)} --- ``electricity travels
like a γ-ray.''

Forward predictions (falsifiers, no fitting): for HeNe lines,
\(\chi(632.8\,\text{nm}) = 89.892^\circ\) and
\(\chi(532\,\text{nm}) = 89.825^\circ\) (modules
\texttt{vp\_light\_angle.py}, \texttt{vp\_electromagnetism.py}).

The earlier Word-document formulation \(\lambda = d/\cos\theta\) with
\(d = 5000\) fm is \textbf{superseded}: it has no \(m\)-chain and is
undefined for \(\lambda < d\) (γ-rays, X-rays). The refined law fixes
\(d\) to the \emph{derived} value \(D = 4852.6\) fm and covers the whole
spectrum via the \(m\)-chain.

\hypertarget{em.7-blackbody-radiation-f}{%
\subsection{\texorpdfstring{EM.7 Blackbody radiation
\textbf{{[}F{]}}}{EM.7 Blackbody radiation {[}F{]}}}\label{em.7-blackbody-radiation-f}}

The Planck spectrum is the product of two lattice facts:
\[ u(\nu,T) \;\propto\; \underbrace{\nu^2}_{\text{geometric mode density}}\;\times\;
   \underbrace{\frac{h\nu}{e^{h\nu/k_BT}-1}}_{\text{rigid-shell high-}\nu\text{ cutoff}}. \]
The first factor alone is Rayleigh--Jeans and diverges (the UV
catastrophe); the rigid shell blocks high-frequency deformation when the
shell energy exceeds the thermal excitation, supplying the Bose cutoff.
Module \texttt{vp\_electromagnetism.py} tabulates \(u_{\mathrm{RJ}}\)
versus \(u_{\mathrm{Planck}}\) (agreement at low \(\nu\), strong
suppression at high \(\nu\)) and solves the Wien transcendental
\(x = 5(1-e^{-x})\) to recover
\(\lambda_{\text{peak}}T = 2.89777\times 10^{-3}\) m·K to
\(6\times 10^{-9}\)\%. A 300 K body peaks at \(\sim 9.7\,\mu\)m
(infrared) with \(\chi \approx 89.95^\circ\) --- transverse light,
consistent with EM.6.

\hypertarget{em.8-refraction-and-dispersion-the-bridge-to-chemistry-fcal}{%
\subsection{\texorpdfstring{EM.8 Refraction and dispersion --- the
bridge to chemistry
\textbf{{[}F{]}/{[}CAL{]}}}{EM.8 Refraction and dispersion --- the bridge to chemistry {[}F{]}/{[}CAL{]}}}\label{em.8-refraction-and-dispersion-the-bridge-to-chemistry-fcal}}

Once light exists, its interaction with matter is geometric. Interfacial
transverse-oscillation matching gives Snell's law
\(n_1\sin\theta_1 = n_2\sin\theta_2\) \textbf{{[}F{]}}; the primary
rainbow angle is \(42.1^\circ\) from Snell + droplet sphere geometry
(Descartes) \textbf{{[}F{]}}; total internal reflection has
\(\sin\theta_c = 1/n\) \textbf{{[}F{]}} (water \(48.6^\circ\), diamond
\(24.4^\circ\)). The absolute refractive index \(n\) is a material
constant \textbf{{[}CAL{]}}. Wavelength-dependent \(n(\lambda)\) orders
the rainbow colors (blue refracts more). These feed directly into the
chemistry chapter's color of coordination compounds (\(d\)--\(d\)
transitions, \(\Delta_{\text{tet}}/\Delta_{\text{oct}} = 4/9\)) and
spectroscopy (modules \texttt{vp\_refraction.py},
\texttt{vp\_light\_angle.py}, \texttt{vp\_crystal\_field.py}).

\hypertarget{em.9-master-ledger-forced-measured-open}{%
\subsection{EM.9 Master ledger --- forced / measured /
open}\label{em.9-master-ledger-forced-measured-open}}

\begin{longtable}[]{@{}
  >{\raggedright\arraybackslash}p{(\columnwidth - 4\tabcolsep) * \real{0.3333}}
  >{\raggedright\arraybackslash}p{(\columnwidth - 4\tabcolsep) * \real{0.3333}}
  >{\raggedright\arraybackslash}p{(\columnwidth - 4\tabcolsep) * \real{0.3333}}@{}}
\toprule\noalign{}
\begin{minipage}[b]{\linewidth}\raggedright
Result
\end{minipage} & \begin{minipage}[b]{\linewidth}\raggedright
Grade
\end{minipage} & \begin{minipage}[b]{\linewidth}\raggedright
Basis / module
\end{minipage} \\
\midrule\noalign{}
\endhead
\bottomrule\noalign{}
\endlastfoot
\(c^2 = B/\rho\) (single surviving longitudinal speed) & {[}F{]}/{[}V{]}
& isostatic \(z{=}2d\); \texttt{vp\_light\_emergence.py} \\
Linear dispersion \(\omega = cq\) & {[}F{]}/{[}V{]} & long-\(\lambda\)
limit; \texttt{vp\_light\_emergence.py} \\
Quantum diameter \(D = 2\lambda_{C,e} = 6\pi^6 r_p\) & {[}F{]} &
wavelength relation; (−19 ppm residual) \\
\(1/r^2\) from \(4\pi r^2\) flux dilution & {[}F{]} & Gauss/Poisson;
\texttt{vp\_electromagnetism.py} \\
Coulomb constant \(k_e = \alpha_{\mathrm{em}}\hbar c/e^2\) &
{[}F{]}+{[}CAL{]} & reproduces \(k_e\) to \(3\times10^{-10}\)\% \\
\(|\mathbf B| = (v/c)|\mathbf E|\), \(\nabla\!\cdot\mathbf B = 0\) &
{[}F{]} & twist geometry; \texttt{vp\_electromagnetism.py} \\
Circular polarization \(=\) rotation (\(|S_3|=S_0\)) & {[}F{]} & Stokes;
\texttt{vp\_electromagnetism.py} \\
Propagation angle \(\sin\chi = \lambda/(mD)\) & {[}F{]}/{[}VP-pred{]} &
\texttt{vp\_light\_angle.py} \\
Blackbody \(=\) mode density \(\times\) rigid cutoff; Wien & {[}F{]} &
\texttt{vp\_electromagnetism.py} \\
Snell, rainbow \(42.1^\circ\), critical angle & {[}F{]} &
\texttt{vp\_refraction.py} \\
Fine-structure coupling \(\alpha_{\mathrm{em}}\) & {[}CAL{]} & measured
input, \textbf{not derived} \\
Long-range / unscreened (global \(U(1)\) Goldstone) & {[}H{]} & symmetry
argument \\
Magnetic / radiative \emph{absolute} vector sector; γγ genesis dynamics
& {[}O{]} & open (physics §14.0.6) \\
\end{longtable}

\hypertarget{em.10-falsification-criteria-f}{%
\subsection{\texorpdfstring{EM.10 Falsification criteria
\textbf{{[}F{]}}}{EM.10 Falsification criteria {[}F{]}}}\label{em.10-falsification-criteria-f}}

Because the forced results have no free parameters, they are refuted if
measurement disagrees:

\begin{itemize}
\tightlist
\item
  If the visible-band propagation angle is \emph{exactly} \(90^\circ\)
  (not \(\approx 89.9^\circ\)), the rotating-quantum / \(m\)-chain
  picture is wrong.
\item
  If \(\chi(632.8\,\text{nm}) \neq 89.892^\circ\) or
  \(\chi(532\,\text{nm}) \neq 89.825^\circ\), the angle law is refuted
  (forward prediction, length → angle, no fitting).
\item
  If the Coulomb field deviates from \(1/r^2\) in three dimensions, the
  flux-dilution origin fails.
\item
  If circular light is not pure rotation (\(|S_3| \neq S_0\)),
  helicity-as-rotation is wrong.
\item
  If the blackbody peak does not satisfy the Wien transcendental, the
  mode-density × rigid-cutoff factorization fails.
\end{itemize}

\hypertarget{em.11-reproducibility}{%
\subsection{EM.11 Reproducibility}\label{em.11-reproducibility}}

All numbers above are produced by deterministic, standard-library Python
modules with self-checks (\texttt{assert}). Determinism is verified by
running each module twice and comparing the stdout sha256.

\begin{verbatim}
python3 vp_light_emergence.py      # light emergence: c²=B/ρ, dispersion, single speed, angle
python3 vp_electromagnetism.py     # charge, 1/r², E/B, polarization, conduction/radiation, blackbody
python3 verify_chemistry.py --root <package root>   # runs all modules + ledger integrity gate
\end{verbatim}

The harness reports execution, determinism, standard-library-only
dependency, and case-ledger integrity across all modules (currently
32/32 PASS with the two EM/light modules added).

\emph{This chapter completes the electromagnetic foundation on which the
chemistry chapter rests: light emerges from the jammed lattice (EM.2),
charge and the Coulomb \(1/r^2\) follow (EM.3), the E/B geometry and
polarization are rotation (EM.4--5), the same wave spans conduction and
radiation by angle (EM.6), and the optical consequences (blackbody,
refraction, color) bridge into chemistry (EM.7--8). It is deliberately
more detailed and fully reproducible relative to the original Korean
working document, and every forced claim is falsifiable.}

\hypertarget{em.12-dynamic-field-equations-poynting-and-causality-addendum-fv}{%
\section{\texorpdfstring{EM.12 Dynamic field equations, Poynting, and
causality (addendum)
\textbf{{[}F{]}/{[}V{]}}}{EM.12 Dynamic field equations, Poynting, and causality (addendum) {[}F{]}/{[}V{]}}}\label{em.12-dynamic-field-equations-poynting-and-causality-addendum-fv}}

\emph{Added 2026-06-14. Closes the dynamic-Maxwell gap flagged in the EM
completeness review (Faraday/Ampère/Poynting were absent from
EM.0--EM.11). Source: \textbf{AQD --- Axiomatic Quantized Dynamics,
NOCAL v1.1, DOI 10.5281/zenodo.17423870}. Reproduced by
\texttt{vp\_em\_field\_dynamics.py}.}

\begin{quote}
The original EM chapter derived the \textbf{static} Coulomb sector
(\(1/r^2\), EM.3), the algebraic ratio \(|\mathbf B|=(v/c)|\mathbf E|\)
and \(\nabla\!\cdot\mathbf B=0\) (EM.4) --- i.e.~\textbf{two of the four
Maxwell statements} (the two divergence facts). It did \textbf{not}
state the \textbf{dynamic} field equations, the Poynting energy flow, or
a causality theorem. This addendum supplies them in the
\textbf{longitudinal (E-sector, pressure \(P\)) form} from the AQD
axioms, and marks the remaining vector-curl sector honestly as open.
\end{quote}

\hypertarget{em.12.1-the-first-order-dynamic-system-f}{%
\subsection{\texorpdfstring{EM.12.1 The first-order dynamic system
\textbf{{[}F{]}}}{EM.12.1 The first-order dynamic system {[}F{]}}}\label{em.12.1-the-first-order-dynamic-system-f}}

The jammed-lattice field obeys a coupled first-order system (AQD axiom
A1) --- pressure \(P\) (the compression = deficit = longitudinal/Coulomb
sector) and velocity \(\mathbf v\): \[
\partial_t P + K\,\nabla\!\cdot\mathbf v = 0,
\qquad
\rho\,\partial_t \mathbf v + \nabla P = 0.
\] Eliminating \(\mathbf v\) gives the wave equation and the speed
already used in EM.2: \[
\Box_c P = 0,\qquad c^2 = \frac{K}{\rho}\;(=B/\rho).
\] \textbf{These two coupled first-order equations are the longitudinal
analogue of Maxwell's curl pair}
(\(\nabla\times\mathbf E=-\partial_t\mathbf B\),
\(\nabla\times\mathbf B=\mu_0\varepsilon_0\partial_t\mathbf E\)): the
time-derivative of each field is sourced by a spatial derivative of the
other, and their combination is the wave. Module
\texttt{vp\_em\_field\_dynamics.py} integrates this system (staggered
leapfrog, deterministic) and recovers the pulse speed
\(c_{\text{measured}}=1.0001\) vs \(c=1\) (\textbf{0.01 \%}).

\hypertarget{em.12.2-poyntings-theorem-and-energy-conservation-fv}{%
\subsection{\texorpdfstring{EM.12.2 Poynting's theorem and energy
conservation
\textbf{{[}F{]}/{[}V{]}}}{EM.12.2 Poynting's theorem and energy conservation {[}F{]}/{[}V{]}}}\label{em.12.2-poyntings-theorem-and-energy-conservation-fv}}

Define the energy density and flux (AQD; equivalent to the second-order
form \(e=\tfrac12[(\partial_tP)^2/c^2+|\nabla P|^2]\),
\(\mathbf S=-(\partial_tP)\nabla P\)): \[
e = \tfrac12\!\left[\frac{P^2}{K} + \rho\,|\mathbf v|^2\right],
\qquad
\mathbf S = P\,\mathbf v,
\qquad
\boxed{\;\partial_t e + \nabla\!\cdot\mathbf S = 0\;}
\] --- the \textbf{local energy-conservation (Poynting) law}, the
longitudinal counterpart of EM's
\(\partial_t u + \nabla\!\cdot(\mathbf E\times\mathbf H)=0\).
Integrating over a lossless domain (periodic, Neumann, time-invariant
Dirichlet, or radiationless infinity) gives
\(\oint\mathbf S\!\cdot\!\mathbf n=0\) and hence \[
\frac{dE}{dt}=0\qquad(\text{AQD QP-0002-001, fully proven}).
\] The module confirms \(\partial_t e+\partial_x S=0\) pointwise (RMS
residual \(\sim 3\times10^{-2}\), discretization scale) and energy
conservation to \(3\times10^{-4}\) over the run. A companion uniqueness
theorem (energy method, QP-0002-002) shows identical initial data give a
unique solution.

\hypertarget{em.12.3-causality-finite-propagation-and-the-light-cone-f}{%
\subsection{\texorpdfstring{EM.12.3 Causality: finite propagation and
the light cone
\textbf{{[}F{]}}}{EM.12.3 Causality: finite propagation and the light cone {[}F{]}}}\label{em.12.3-causality-finite-propagation-and-the-light-cone-f}}

The same energy structure forces a \textbf{hard light cone} (AQD
theorems F-A/B/C): \[
u_{\max}=c,\qquad \operatorname{supp} e(\tau,\cdot)\ \subseteq\ \operatorname{supp} e(0,\cdot)\,\oplus\,B_{c\tau}.
\] Energy and information cannot reach beyond \(c\tau\) of the initial
support. The module verifies this directly: the pulse support stays
inside the cone (expansion \(0.00 \le c\tau\)). At the quantum level the
AQD upgrades this to \textbf{microcausality} --- with
\(\Phi:=P/\sqrt K\), \([\Phi(x),\Phi(y)]=0\) for spacelike \((x-y)\) via
the Pauli--Jordan function \(D(x)\) whose support lies on the light cone
(the cone speed is fixed at \(c\) by \(c^2=K/\rho\)), and the
arrival-time bound \(t_{\text{arr}}\ge \text{dist}/c\) follows from
commutator support. No calibration enters.

\hypertarget{em.12.4-what-this-closes-and-what-remains-open-o}{%
\subsection{\texorpdfstring{EM.12.4 What this closes, and what remains
open
\textbf{{[}O{]}}}{EM.12.4 What this closes, and what remains open {[}O{]}}}\label{em.12.4-what-this-closes-and-what-remains-open-o}}

\textbf{Closed (this addendum):} the \emph{dynamic} field equations, the
Poynting energy flow and its exact conservation, uniqueness, and
causality (classical light cone + quantum microcausality). Together with
EM.2 (wave/\(c^2=B/\rho\)) and EM.3 (static \(1/r^2\)), the
\textbf{longitudinal (E-sector)} electrodynamics is now complete and
reproducible.

\textbf{Still open --- honestly marked {[}O{]} (unchanged from EM.9):}
the \textbf{full vector \(\mathbf E/\mathbf B\) curl structure} ---
Faraday's \(\nabla\times\mathbf E=-\partial_t\mathbf B\) and
Ampère--Maxwell's \(\nabla\times\mathbf B\) in \emph{vector} form, with
the \textbf{magnetic field as an independent transverse pseudovector}
and its \emph{absolute} magnitude. The AQD treatment is
scalar/longitudinal (\(P\), \(\mathbf v\) with \(\mathbf v\)
irrotational); it yields the compression (Coulomb/\(E\)) sector
rigorously but not the independent transverse-vector
(magnetic/radiative) sector. This is exactly the EM.9 line
\emph{``magnetic/radiative absolute vector sector --- {[}O{]} (physics
§14.0.6)''}. It is \textbf{not} resolved here; it remains the principal
open structural item of the EM chapter.

\hypertarget{em.12.5-updated-master-ledger-rows}{%
\subsection{EM.12.5 Updated master-ledger
rows}\label{em.12.5-updated-master-ledger-rows}}

\begin{longtable}[]{@{}
  >{\raggedright\arraybackslash}p{(\columnwidth - 4\tabcolsep) * \real{0.3333}}
  >{\raggedright\arraybackslash}p{(\columnwidth - 4\tabcolsep) * \real{0.3333}}
  >{\raggedright\arraybackslash}p{(\columnwidth - 4\tabcolsep) * \real{0.3333}}@{}}
\toprule\noalign{}
\begin{minipage}[b]{\linewidth}\raggedright
Result
\end{minipage} & \begin{minipage}[b]{\linewidth}\raggedright
Grade
\end{minipage} & \begin{minipage}[b]{\linewidth}\raggedright
Basis / module
\end{minipage} \\
\midrule\noalign{}
\endhead
\bottomrule\noalign{}
\endlastfoot
First-order system \(\partial_tP+K\nabla\!\cdot\mathbf v=0\),
\(\rho\partial_t\mathbf v+\nabla P=0\) & {[}F{]} & AQD A1;
\texttt{vp\_em\_field\_dynamics.py} \\
Poynting law \(\partial_t e+\nabla\!\cdot\mathbf S=0\),
\(\mathbf S=P\mathbf v\) & {[}F{]}/{[}V{]} & AQD QP-0002-001; residual
\$\sim\$3e-2 \\
Energy conservation \(dE/dt=0\) (lossless) + uniqueness & {[}F{]} & AQD
QP-0002-001/002; dev 3e-4 \\
Causality: light cone, \(u_{\max}=c\); microcausality & {[}F{]} & AQD
F-A/B/C, QP-0002-U121 \\
Full vector \(\mathbf E/\mathbf B\) curl (Faraday/Ampère, magnetic
absolute) & \textbf{{[}O{]}} & open --- scalar/longitudinal AQD does not
supply it \\
\end{longtable}

\emph{This addendum upgrades the EM chapter from ``static sector +
algebraic \(B/E\) ratio'' to ``\textbf{complete longitudinal
electrodynamics with Poynting flow and causality},'' using the AQD
axiomatic field dynamics. The dynamic-Maxwell gap is closed in the
E-sector; the independent vector (magnetic/radiative) sector remains the
chapter's honestly-declared open item.}

\hypertarget{chemistry-from-a-single-anchor-the-jamming-account}{%
\section{Chemistry from a Single Anchor: The Jamming
Account}\label{chemistry-from-a-single-anchor-the-jamming-account}}

\textbf{VP Chemistry \& Electromagnetism Whitepaper --- Chapter CH} DOI
(reserved): 10.5281/zenodo.20680541 · Companion to VP Theory (physics)
and to Chapter EM Status grades: \textbf{{[}F{]}} forced (zero free
parameters) · \textbf{{[}F?{]}} forced-candidate (model-dependent or
\textless1.5\%) · \textbf{{[}CAL{]}} calibration input (measured
constant, used consistently) · \textbf{{[}V{]}} simulation-measured ·
\textbf{{[}O{]}} open (honestly unsolved). Every numeric claim is
reproduced by a deterministic, standard-library module (2× sha256
identical); module and case-ledger anchors are given inline. Verified
numbers quoted below are the actual module outputs.

\begin{quote}
Everything starts locked, is derived, and is settled by a verdict. One
table separates what is \emph{forced} from what is \emph{measured input}
from what is \emph{open}. Chemistry unfolds from the electron (a
rotational dent) using only \(\pi\)-geometry and the electromagnetic
\(1/r^2\) --- the same substrate from which light emerges (Chapter EM).
\end{quote}

\hypertarget{ch.0-governance-and-grading-inherited-from-vp-theory-1}{%
\subsection{CH.0 Governance and grading (inherited from VP Theory
§1)}\label{ch.0-governance-and-grading-inherited-from-vp-theory-1}}

This chapter inherits No-Tuning / LOCK / Gate. \textbf{No-Tuning:}
changing an input after seeing a result to hit a target is forbidden.
The clean final logic is presented here; the trial-and-error of
development lives, reproducibly, in the modules and the case ledger.
Forced results \textbf{{[}F{]}} have no adjustable parameter --- they
cannot be tuned, so they are \emph{refuted} if they disagree with
measurement.

\hypertarget{ch.1-the-anti-tuning-chain-the-central-defense}{%
\subsection{CH.1 The anti-tuning chain --- the central
defense}\label{ch.1-the-anti-tuning-chain-the-central-defense}}

A trained reviewer immediately suspects parameter-fitting. The rebuttal
has three prongs.

\textbf{Prong A --- forced results are dimensionless geometry with
nothing to adjust.} The tetrahedral angle \(\arccos(-1/3)\), the ratio
\(\Delta_{\text{tet}}/\Delta_{\text{oct}} = 4/9\), the FCC packing
fraction \(\pi/(3\sqrt2)\), the shell-filling integers
\([2,8,8,18,18,32]\) --- these are pure geometry/integers with no knob.

\textbf{Prong B --- calibration inputs are used as forward predictions.}
Measured constants (spectroscopic \(B,\nu\); bond enthalpy \(\Delta H\);
\(\mathrm pK_a\); standard potential \(E^\circ\)) are each fixed
\emph{once} and used consistently. Quantities \emph{derived} from them
--- standard entropy, equilibrium constants, cell potentials --- then
match measurement; that is prediction, not fitting.

\textbf{Prong C --- honest open items prove no tuning.} A tuner would
have ``solved'' absolute multi-electron ionization energies; instead the
failure is recorded (mean \textasciitilde22\% for main group, electron
affinity wrong in sign 8/18 by the Slater method) and marked
\textbf{{[}O{]}}.

\hypertarget{ch.2-foundation-anchor-atomic-scale-f}{%
\subsection{\texorpdfstring{CH.2 Foundation: anchor → atomic scale
\textbf{{[}F{]}}}{CH.2 Foundation: anchor → atomic scale {[}F{]}}}\label{ch.2-foundation-anchor-atomic-scale-f}}

The only observational input is the electron mass \(m_e\) (the anchor).
With the EM coupling \(\alpha_{\mathrm{em}}\) \textbf{{[}CAL{]}} and
\(\pi\) (using the \emph{reduced} Compton wavelength
\(\bar\lambda_{C,e} = \hbar/(m_e c) = \lambda_{C,e}/2\pi\)):
\[ a_0 = \frac{\bar\lambda_{C,e}}{\alpha_{\mathrm{em}}} = \frac{\hbar}{m_e c\,\alpha_{\mathrm{em}}},
   \qquad \mathrm{Ry} = \frac{\alpha_{\mathrm{em}}^2 m_e c^2}{2}. \]
Module \texttt{vp\_atomic.py} gives \(a_0 = 0.529177\) Å (Δ +0.0000\%)
and \(\mathrm{Ry} = 13.6057\) eV (Δ −0.0001\%); hydrogenic
\(\mathrm{IE} = \mathrm{Ry}\,Z^2\) to \textless0.06\%. Shell structure
follows from the Aufbau degeneracy \(2(2l+1)\), giving noble-gas
closures \([2,10,18,36,54,86]\) and period lengths \([2,8,8,18,18,32]\)
exactly \textbf{{[}F{]}}. Absolute many-electron IE is only approximate
(screening) and is \textbf{{[}O{]}}.

\hypertarget{ch.3-periodic-structure-valence-and-oxidation-ff}{%
\subsection{\texorpdfstring{CH.3 Periodic structure: valence and
oxidation
\textbf{{[}F{]}/{[}F?{]}}}{CH.3 Periodic structure: valence and oxidation {[}F{]}/{[}F?{]}}}\label{ch.3-periodic-structure-valence-and-oxidation-ff}}

Valence equals the number of outer-shell electrons (1--8, main group)
\textbf{{[}F{]}}. The octet rule fixes ionic-compound stoichiometry:
module \texttt{vp\_valence.py} predicts 9/9 formulas (NaCl, MgCl₂,
Al₂O₃, CaO, \ldots). Periodic trends (radius ↓, electronegativity ↑
across a period; radius ↑ down a group) are \textbf{{[}F?{]}}.
Transition metals have valence \(4s+3d\) with variable oxidation states
(Mn +2\ldots+7) from the partially filled \(d\)-shell; the
half/fully-filled exceptions Cr (\(3d^54s^1\)), Cu (\(3d^{10}4s^1\))
follow from sub-shell stability \textbf{{[}F{]}}.

\hypertarget{ch.4-chemical-bonding-and-vsepr-geometry-f}{%
\subsection{\texorpdfstring{CH.4 Chemical bonding and VSEPR geometry
\textbf{{[}F{]}}}{CH.4 Chemical bonding and VSEPR geometry {[}F{]}}}\label{ch.4-chemical-bonding-and-vsepr-geometry-f}}

Geometry is the electromagnetic \(1/r^2\) repulsion of electron domains
on a sphere (the Thomson problem). Module
\texttt{vp\_molecular\_geometry.py} \emph{computes} (seed-independent)
the angles:
\[ \text{linear } 180^\circ,\quad \text{trigonal } 120^\circ,\quad
   \boxed{\text{tetrahedral } \arccos(-\tfrac13) = 109.47^\circ},\quad \text{octahedral } 90^\circ. \]
All have zero free parameters \textbf{{[}F{]}}; the simulation
reproduces the closed form \(\arccos(-1/3)=109.47^\circ\). Lone-pair
compression (CH₄ → NH₃ → H₂O:
\(109.5^\circ \to 107^\circ \to 104.5^\circ\)) has the correct direction
\textbf{{[}F?{]}}; the precise magnitude is \textbf{{[}O{]}}. Bond
rupture follows a scale-invariant \(\sqrt2\) diagonal-stretch geometry
shared with nuclear fission (Chapter EM links the same \(1/r^2\));
absolute bond energy needs per-bond anharmonicity and is
\textbf{{[}O{]}}.

\hypertarget{ch.5-coordination-chemistry-color-and-magnetism-f}{%
\subsection{\texorpdfstring{CH.5 Coordination chemistry: color and
magnetism
\textbf{{[}F{]}}}{CH.5 Coordination chemistry: color and magnetism {[}F{]}}}\label{ch.5-coordination-chemistry-color-and-magnetism-f}}

Ligand point charges (EM \(1/r^2\)) split the geometric \(d\)-orbitals
(spherical harmonics). Module \texttt{vp\_crystal\_field.py} gives the
octahedral \(e_g\)↑ / \(t_{2g}\)↓ pattern (tetrahedral inverted) and the
analytic ratio
\[ \frac{\Delta_{\text{tet}}}{\Delta_{\text{oct}}} = \frac{4}{9} = 0.4444 \quad\text{(exact)}. \]
High/low-spin unpaired-electron counts give magnetic moments (Fe³⁺ 5/1
unpaired) \textbf{{[}F{]}}. Color is the \(d\)--\(d\) transition
\(\Delta \to \lambda\) (interpreted optically via Chapter EM);
multi-band precise color is \textbf{{[}O{]}}, the spectrochemical series
is \textbf{{[}CAL{]}}.

\hypertarget{ch.6-light-optics-color-see-chapter-em-fcal}{%
\subsection{\texorpdfstring{CH.6 Light, optics, color --- see Chapter EM
\textbf{{[}F{]}/{[}CAL{]}}}{CH.6 Light, optics, color --- see Chapter EM {[}F{]}/{[}CAL{]}}}\label{ch.6-light-optics-color-see-chapter-em-fcal}}

The optical consequences (propagation angle \(\sin\chi = \lambda/(mD)\),
Snell, rainbow \(42.1^\circ\), total-internal-reflection critical angle,
blackbody, \(d\)--\(d\) color) are developed in \textbf{Chapter EM}.
Chemistry \emph{consumes} the electromagnetic mechanism from physics
§10/§14; it does not re-derive it. The boundary is explicit: fundamental
EM (light as rotational transverse wave, \(c^2=B/\rho\), the \(1/r^2\)
Green function) is physics; the chemical/optical consequences are here.

\hypertarget{ch.7-thermodynamics-rotational-energy-temperature-ff}{%
\subsection{\texorpdfstring{CH.7 Thermodynamics: rotational energy →
temperature
\textbf{{[}F{]}/{[}F?{]}}}{CH.7 Thermodynamics: rotational energy → temperature {[}F{]}/{[}F?{]}}}\label{ch.7-thermodynamics-rotational-energy-temperature-ff}}

Temperature is a rotational-energy scale; rotational levels
\(E_J = B\,hc\,J(J{+}1)\). Heat-capacity degrees-of-freedom: monatomic
\(\tfrac32 R\), diatomic \(\tfrac52 R\), solid Dulong--Petit \(3R\)
(\textless3\%). Module \texttt{vp\_statistical\_thermo.py} reproduces
the H₂ heat-capacity \textbf{staircase} (\(\tfrac32 R\) at 20 K →
\(\tfrac52 R\) at 298 K → \(\tfrac72 R\) above \textasciitilde6332 K)
with activation temperatures \(T_{\text{rot}} = 88\) K,
\(T_{\text{vib}} = 6332\) K from \(B,\nu\); \(C_v\)(H₂) 20.81 vs 20.4
(+2.0\%) \textbf{{[}F{]}}. Standard molar entropy (Sackur--Tetrode +
rotation + vibration + electronic) is \textless0.3\% for H₂·N₂·CO·O₂; O₂
includes the known triplet (paramagnetic) \(R\ln 3\) \textbf{{[}F?{]}}.
The spectroscopic constants \(B,\nu\) are per-molecule
\textbf{{[}CAL{]}}.

\hypertarget{ch.8-chemical-equilibrium-f}{%
\subsection{\texorpdfstring{CH.8 Chemical equilibrium
\textbf{{[}F{]}}}{CH.8 Chemical equilibrium {[}F{]}}}\label{ch.8-chemical-equilibrium-f}}

Direction is set by \(\Delta G = \Delta H - T\Delta S\); dissociation
occurs when \(T\Delta S > \Delta H\). The dissociation temperature is
\[ T^* = \frac{\Delta H}{\Delta S}. \] Module
\texttt{vp\_equilibrium.py} gives \(T^*\): N₂ 8214 K \textgreater{} H₂
4400 K \textgreater{} O₂ 4210 K \textgreater{} Cl₂ 2272 K --- the
\textbf{bond-strength order} \textbf{{[}F{]}}. Reaction entropies from
partition functions are \textless0.4\%. The van 't Hoff \(K(T)\) (H₂ ⇌
2H, with \(K=1\) at \(T^*\)) is \textbf{{[}F?{]}} (\(\Delta H\) is
\textbf{{[}CAL{]}}). Le Chatelier directions (dissociation
\(\Delta S>0\) high-T driven; Haber \(\Delta S<0\) high-T reversed)
\textbf{{[}F{]}}.

\hypertarget{ch.9-reaction-kinetics-f}{%
\subsection{\texorpdfstring{CH.9 Reaction kinetics
\textbf{{[}F{]}}}{CH.9 Reaction kinetics {[}F{]}}}\label{ch.9-reaction-kinetics-f}}

The activation energy is the threshold to reach the transition state (a
\(\sqrt2\) bond-stretch); the fraction crossing is the Boltzmann tail
\(\exp(-E_a/RT)\). Module \texttt{vp\_kinetics.py} reproduces the
Arrhenius temperature dependence, the room-temperature \textbf{10 K
doubling rule} (\(E_a\approx 50\) kJ/mol gives ×1.9 per 10 K), and
recovery of \(E_a\) from the Arrhenius slope \(-E_a/R\). Mean molecular
speed \(\sqrt{8k_BT/\pi m}\) (N₂ 298 K → 475 m/s) and the Eyring
universal factor \(k_BT/h = 6.2\times10^{12}\) s⁻¹ follow
\textbf{{[}F{]}}. Absolute \(E_a\) (transition-state structure) and
tunneling are \textbf{{[}O{]}}.

\hypertarget{ch.10-solution-chemistry-proton-transfer-equilibrium-f}{%
\subsection{\texorpdfstring{CH.10 Solution chemistry: proton-transfer
equilibrium
\textbf{{[}F{]}}}{CH.10 Solution chemistry: proton-transfer equilibrium {[}F{]}}}\label{ch.10-solution-chemistry-proton-transfer-equilibrium-f}}

With \(K_w = 10^{-14}\), pure water has pH 7. The weak-acid pH solves
\(K_a = x^2/(C-x)\); module \texttt{vp\_solution\_chemistry.py} gives
acetic acid 0.1 M → pH 2.88 (measured 2.87). Henderson--Hasselbalch
buffering pH \(= \mathrm pK_a + \log([A^-]/[HA])\) and titration
equivalence-point jumps follow \textbf{{[}F{]}}. Individual
\(\mathrm pK_a\) are \textbf{{[}CAL{]}}; absolute \(\mathrm pK_a\)
prediction (solvation) is \textbf{{[}O{]}}.

\hypertarget{ch.11-electrochemistry-electron-rotational-dent-transfer-f}{%
\subsection{\texorpdfstring{CH.11 Electrochemistry: electron (rotational
dent) transfer
\textbf{{[}F{]}}}{CH.11 Electrochemistry: electron (rotational dent) transfer {[}F{]}}}\label{ch.11-electrochemistry-electron-rotational-dent-transfer-f}}

Charge is rotation (Chapter EM); redox is rotational-dent transfer,
symmetric with proton transfer (CH.10). The cell potential is
\[ E^\circ_{\text{cell}} = E^\circ_{\text{cathode}} - E^\circ_{\text{anode}}, \qquad
   \Delta G^\circ = -nFE^\circ,\qquad E = E^\circ - \frac{RT}{nF}\ln Q. \]
Module \texttt{vp\_electrochemistry.py} gives the Daniell cell
\(E^\circ_{\text{cell}} = 1.10\) V (matches measured) and
\(\Delta G^\circ = -212\) kJ/mol \textbf{{[}F{]}}. The
standard-potential series (Li −3.04 → F₂ +2.87) links to IE and
electronegativity; individual values are \textbf{{[}CAL{]}}.

\hypertarget{ch.12-organic-chemistry-carbon-skeleton-on-em-geometry-f}{%
\subsection{\texorpdfstring{CH.12 Organic chemistry: carbon skeleton on
EM geometry
\textbf{{[}F{]}}}{CH.12 Organic chemistry: carbon skeleton on EM geometry {[}F{]}}}\label{ch.12-organic-chemistry-carbon-skeleton-on-em-geometry-f}}

Carbon's tetravalence plus EM repulsion gives hybridization geometry
(the same VSEPR engine as CH.4): module
\texttt{vp\_organic\_chemistry.py} gives sp³ tetrahedral
\(109.5^\circ\), sp² planar \(120^\circ\), sp linear \(180^\circ\). Bond
order--length--strength: \[
\begin{array}{lccc}
 & \text{order} & \text{length} & \text{strength (kJ/mol)}\\
\text{C–C} & 1\ (1\sigma) & 154\ \text{pm} & 348\\
\text{C=C} & 2\ (1\sigma+1\pi) & 134\ \text{pm} & 614\\
\text{C≡C} & 3\ (1\sigma+2\pi) & 120\ \text{pm} & 839
\end{array}
\] The degree of unsaturation \((2C+2+N-H-X)/2\) gives benzene 4
\textbf{{[}F{]}}. Functional-group reactivity (acid--base, redox
applied) is \textbf{{[}F?{]}}; absolute rates and stereoselectivity are
\textbf{{[}O{]}}.

\hypertarget{ch.13-materials-and-nuclear-summary-ff}{%
\subsection{\texorpdfstring{CH.13 Materials and nuclear (summary)
\textbf{{[}F{]}/{[}F?{]}}}{CH.13 Materials and nuclear (summary) {[}F{]}/{[}F?{]}}}\label{ch.13-materials-and-nuclear-summary-ff}}

Crystal packing fractions are pure geometry: FCC/HCP
\(\pi/(3\sqrt2) = 0.7405\), BCC 0.6802, SC 0.5236 (module
\texttt{vp\_crystal\_packing.py}); FCC metal densities \textless2\%
\textbf{{[}F{]}/{[}F?{]}}. Nuclear magic-number skeleton from
Gauss-lattice point counts \(N(R^2)+1\): module
\texttt{vp\_gauss\_shells.py} gives \textbf{2, 8, 20, 28, 82} (5/7
parameter-free); 50, 126 require spin--orbit \textbf{{[}CAL{]}}. SEMF
fission/fusion gives the iron peak \(A=58\) and the critical \(Z^2/A\)
\textbf{{[}F?{]}}. Theoretical bounds (\(\sigma_{th}\), \(\gamma_s\))
are bounds, not predictions --- honestly \textbf{{[}F{]}+{[}O{]}}.

\hypertarget{ch.14-master-ledger-forced-measured-input-open}{%
\subsection{CH.14 Master ledger --- forced / measured input /
open}\label{ch.14-master-ledger-forced-measured-input-open}}

\textbf{{[}F{]} forced (zero free parameters --- cannot be tuned):}
shell filling and period lengths; valence, oxidation, ionic formulas;
all VSEPR angles (\(\arccos(-1/3)\) etc.);
\(\Delta_{\text{tet}}/ \Delta_{\text{oct}}=4/9\); packing fractions;
propagation angle \(\sin\chi=\lambda/(mD)\); rainbow \(42^\circ\);
Snell, critical angle; \(C_v(T)\) staircase and activation temperatures;
dissociation- temperature order and Le Chatelier; Arrhenius, 10 K
doubling, \(k_BT/h\); pH, buffers, titration; \(E^\circ_{\text{cell}}\),
\(\Delta G=-nFE^\circ\), Nernst; hybridization, bond order,
unsaturation; Gauss magic-number skeleton.

\textbf{{[}CAL{]} calibration input (measured constants, used
consistently --- never tuned per case):} spectroscopic \(B,\nu\); bond
enthalpy \(\Delta H\); \(\mathrm pK_a\); standard potential \(E^\circ\);
spectrochemical series; absolute refractive index \(n\); spin--orbit
strength (magic 50, 126); the fine-structure coupling
\(\alpha_{\mathrm{em}}\).

\textbf{{[}O{]} open (honest --- not disguised as solved):} absolute
multi-electron IE (screening); electron affinity (Slater sign failure);
absolute \(\mathrm pK_a\), \(E^\circ\) (solvation); multi-band precise
color; absolute bond energy (anharmonicity); \(d\)-block absolute IE;
absolute rates and stereoselectivity; nuclear shell-correction masses;
pure-\(\pi\) geometry of 50, 126.

\hypertarget{ch.15-two-unifications-sphere-geometry-and-rotation}{%
\subsection{CH.15 Two unifications --- sphere geometry and
rotation}\label{ch.15-two-unifications-sphere-geometry-and-rotation}}

\textbf{Sphere geometry:} molecular shape = points \emph{on} a sphere
(Thomson/VSEPR) → \(\arccos(-1/3)\); nuclear shells = lattice points
\emph{inside} a sphere (Gauss) → magic-number skeleton; crystal density
= sphere \emph{packing} (Kepler) → 0.7405; \(d\)-orbital splitting =
spherical harmonics × ligand EM → 4/9.

\textbf{Rotation:} the electron is a rotational dent; charge is rotation
handedness; light is the rotational transverse wave (Chapter EM);
temperature is rotational excitation; equilibrium, rate, and
electrochemistry all ride on rotational energy. The single anchor
(electron = rotational dent) opens, and the last sector
(electrochemistry = rotational-dent transfer) closes the circle.

\hypertarget{ch.16-falsification-criteria-f}{%
\subsection{\texorpdfstring{CH.16 Falsification criteria
\textbf{{[}F{]}}}{CH.16 Falsification criteria {[}F{]}}}\label{ch.16-falsification-criteria-f}}

Forced results cannot be re-fitted, so they are refuted on disagreement:

\begin{itemize}
\tightlist
\item
  If the tetrahedral angle is not \(\arccos(-1/3)=109.47^\circ\), the
  VSEPR-EM derivation is discarded.
\item
  If \(\Delta_{\text{tet}}/\Delta_{\text{oct}} \neq 4/9\), the
  crystal-field geometry is discarded.
\item
  If the rainbow angle is not \(\approx 42^\circ\), the
  light-angle/refraction derivation is discarded.
\item
  If H₂ heat capacity shows no \(\tfrac32\to\tfrac52\to\tfrac72 R\)
  staircase, or the activation temperatures disagree with \(B,\nu\), the
  partition-function account is discarded.
\item
  If dissociation temperatures do not follow the bond-strength order,
  the equilibrium derivation is discarded.
\item
  The committed light-angle predictions
  \(\chi(632.8\,\text{nm})=89.892^\circ\),
  \(\chi(532\,\text{nm})=89.825^\circ\) are falsifiable forward
  predictions (Chapter EM).
\end{itemize}

\hypertarget{ch.17-reproducibility}{%
\subsection{CH.17 Reproducibility}\label{ch.17-reproducibility}}

Thirty deterministic, standard-library modules and a 114-row case ledger
(\texttt{cases\_fixed.csv}) reproduce this logic; the light/EM additions
add a 16-row ledger (\texttt{cases\_em.csv}). Determinism is verified by
running each module twice and comparing the stdout sha256.

\begin{verbatim}
python3 research/foundation/vp_molecular_geometry.py   # tetrahedral 109.47°
python3 research/foundation/vp_crystal_field.py        # 4/9
python3 verify_chemistry.py --root <package root>      # all modules + ledger integrity gate
\end{verbatim}

The harness reports execution, determinism, standard-library-only
dependency, and case-ledger integrity. With the light and
electromagnetism modules included, \textbf{32/32 PASS}.

\emph{Together with Chapter EM, this completes the program: light
emerges from the jammed lattice and the electromagnetic \(1/r^2\)
follows (Chapter EM); on that substrate, atoms, bonds, geometry, color,
heat, equilibrium, rate, solution, electrochemistry, and organic
structure are forced by sphere geometry and rotation (this chapter).
What is forced is falsifiable; what is measured is labelled; what is
open is admitted --- and all of it is reproducible.}

\hypertarget{conduction-in-copper-and-magnetism-in-iron-two-fates-of-the-d-shell}{%
\section{\texorpdfstring{Conduction in Copper and Magnetism in Iron: Two
Fates of the
\emph{d}-Shell}{Conduction in Copper and Magnetism in Iron: Two Fates of the d-Shell}}\label{conduction-in-copper-and-magnetism-in-iron-two-fates-of-the-d-shell}}

\textbf{VP Chemistry \& Electromagnetism Whitepaper --- Chapter CM} DOI
(reserved): 10.5281/zenodo.20680541 · Builds on Chapter EM
(electromagnetism) and Chapter CH (chemistry). Status grades as in
Chapter EM. Numbers reproduced by \texttt{vp\_copper\_conduction.py} and
\texttt{vp\_iron\_magnetism.py} (deterministic, standard-library,
self-checking).

\begin{quote}
Copper conducts and iron magnetizes for the same root reason, read two
ways. Copper's inner \emph{d}-shell is \textbf{full}, freeing its outer
electron into a conducting sea; iron's \emph{d}-shell is \textbf{partly
empty}, leaving unpaired rotations that align into a magnet. A filled
shell makes a conductor; an unfilled shell makes a magnet.
\end{quote}

\hypertarget{cm.0-the-two-questions}{%
\subsection{CM.0 The two questions}\label{cm.0-the-two-questions}}

This chapter answers two concrete questions about metals using the
electromagnetism of Chapter EM:

\begin{enumerate}
\def\labelenumi{\arabic{enumi}.}
\tightlist
\item
  \textbf{Electricity is an electromagnetic wave --- how does it pass
  \emph{through the inside} of copper,} when copper (a metal)
  \emph{reflects} ordinary light?
\item
  \textbf{What is the principle of iron's magnetism} --- and in
  particular, why iron (and cobalt, nickel) but not copper, and not even
  manganese, which has \emph{more} unpaired electrons than iron? (The
  physics whitepaper took the ``why Fe and not Cu'' criterion as an
  empirical input, §14; here it is derived.)
\end{enumerate}

\hypertarget{cm.1-the-paradox}{%
\subsection{CM.1 The paradox}\label{cm.1-the-paradox}}

Light is a transverse electromagnetic wave, and metals reflect it ---
that is why copper is shiny and opaque. Yet electricity is \emph{also}
electromagnetic, and it travels straight through the body of a copper
wire. How can the same medium reflect one electromagnetic disturbance
and transmit another?

The resolution is the propagation angle of Chapter EM (§EM.6).
Electricity is the \textbf{extreme longitudinal limit} (\(\chi \to 0\))
of the electromagnetic wave, while visible light is the
\textbf{transverse limit} (\(\chi \to 90^\circ\)). The same
free-electron sea reflects the transverse wave and transmits the
longitudinal one. Same medium, opposite fates, set by angle.

\hypertarget{cm.2-coppers-free-electron-sea-fcal}{%
\subsection{\texorpdfstring{CM.2 Copper's free-electron sea
\textbf{{[}F{]}/{[}CAL{]}}}{CM.2 Copper's free-electron sea {[}F{]}/{[}CAL{]}}}\label{cm.2-coppers-free-electron-sea-fcal}}

Copper is \([\mathrm{Ar}]\,3d^{10}4s^1\): the inner \(3d\) shell is
\textbf{full}, and the single \(4s\) electron is delocalized into a
conducting sea. The free-electron density is
\[ n = \frac{\rho\,N_A\,Z_{\text{free}}}{M} = 8.49\times 10^{28}\ \mathrm{m^{-3}} \]
(one free electron per atom; \(\rho, M\) measured \textbf{{[}CAL{]}}).
This sea is the ``internal quanta'' through which conduction propagates.

\hypertarget{cm.3-why-transverse-light-is-reflected-the-plasma-frequency-fcal}{%
\subsection{\texorpdfstring{CM.3 Why transverse light is reflected ---
the plasma frequency
\textbf{{[}F{]}/{[}CAL{]}}}{CM.3 Why transverse light is reflected --- the plasma frequency {[}F{]}/{[}CAL{]}}}\label{cm.3-why-transverse-light-is-reflected-the-plasma-frequency-fcal}}

The free-electron sea has a natural oscillation frequency, the plasma
frequency
\[ \omega_p = \sqrt{\frac{n e^2}{\varepsilon_0 m_e}} = 1.64\times 10^{16}\ \mathrm{rad/s}
   \;\;(f_p = 2.6\times 10^{15}\ \mathrm{Hz},\ \lambda_p = 115\ \mathrm{nm},\ \text{ultraviolet}). \]
Below \(\omega_p\) --- radio, infrared, \textbf{all of visible light}
--- an incident transverse wave drives the free electrons, which
re-radiate and \textbf{reflect} it: this is exactly why metals are shiny
and opaque (\texttt{vp\_copper\_conduction.py}, §1). Above \(\omega_p\)
(X-rays) the electrons cannot follow and the metal becomes transparent.
The transverse electromagnetic wave is turned back by the very sea that
carries the current.

\hypertarget{cm.4-why-longitudinal-electricity-is-transmitted-conduction-as-chi-to-0-f}{%
\subsection{\texorpdfstring{CM.4 Why longitudinal electricity is
transmitted --- conduction as \(\chi \to 0\)
\textbf{{[}F{]}}}{CM.4 Why longitudinal electricity is transmitted --- conduction as \textbackslash chi \textbackslash to 0 {[}F{]}}}\label{cm.4-why-longitudinal-electricity-is-transmitted-conduction-as-chi-to-0-f}}

The longitudinal mode is different. A longitudinal (\(\chi \to 0\))
electromagnetic disturbance is a \textbf{density/pressure wave of the
electron sea} that threads straight along the wire --- straighter even
than a \(\gamma\)-ray. It is not reflected because it is not a
transverse oscillation of the surface electrons; it propagates
\emph{through} the internal quanta of the copper atoms, not through
vacuum or air (physics §14). This longitudinal disturbance is the
current. The angle that makes light reflect (transverse) is the angle
that makes electricity transmit (longitudinal). Two physical facts fix
conduction at \(\chi \to 0\): small amplitude (a low-amplitude
disturbance threads straight; if the amplitude grows the path tilts
transverse and the signal \emph{leaves} the copper as radiative loss)
and medium (it runs through the atomic-internal quanta).

\hypertarget{cm.5-the-skin-effect-fcal}{%
\subsection{\texorpdfstring{CM.5 The skin effect
\textbf{{[}F{]}/{[}CAL{]}}}{CM.5 The skin effect {[}F{]}/{[}CAL{]}}}\label{cm.5-the-skin-effect-fcal}}

The transverse and longitudinal pictures are quantitatively distinct. A
transverse wave that does enter the metal decays within the skin depth
\[ \delta = \sqrt{\frac{2}{\mu_0\,\sigma\,\omega}}: \quad
   8.4\ \mathrm{mm}\ (60\ \mathrm{Hz}) \to 2.1\ \mu\mathrm{m}\ (1\ \mathrm{GHz}). \]
High-frequency transverse fields are confined to a micron-thin surface
layer, while the direct/low-frequency longitudinal current fills the
whole cross-section (\texttt{vp\_copper\_conduction.py}, §3).

\hypertarget{cm.6-three-velocities-drift-signal-fermi-fcal}{%
\subsection{\texorpdfstring{CM.6 Three velocities --- drift, signal,
Fermi
\textbf{{[}F{]}/{[}CAL{]}}}{CM.6 Three velocities --- drift, signal, Fermi {[}F{]}/{[}CAL{]}}}\label{cm.6-three-velocities-drift-signal-fermi-fcal}}

A common confusion is dissolved by separating three speeds
(\texttt{vp\_copper\_conduction.py}, §4):

\begin{itemize}
\tightlist
\item
  \textbf{Drift velocity} \(v_d = J/(ne)\): for
  \(J = 10^6\ \mathrm{A/m^2}\), only \(0.074\ \mathrm{mm/s}\) --- the
  electrons themselves crawl.
\item
  \textbf{Signal velocity} \(\sim c\): the longitudinal field
  disturbance travels near light speed, so a lamp lights essentially
  instantly even though the electrons barely move.
\item
  \textbf{Fermi velocity} \(v_F = 1.57\times 10^6\ \mathrm{m/s}\): the
  electrons at the Fermi surface move fast, but their \emph{net} drift
  is the slow \(v_d\).
\end{itemize}

The predicted Fermi energy
\(E_F = \hbar^2(3\pi^2 n)^{2/3}/(2m_e) = 7.04\ \mathrm{eV}\) matches the
measured \(\approx 7.0\ \mathrm{eV}\).

\hypertarget{cm.7-why-copper-conducts-so-well-fcal}{%
\subsection{\texorpdfstring{CM.7 Why copper conducts so well
\textbf{{[}F{]}/{[}CAL{]}}}{CM.7 Why copper conducts so well {[}F{]}/{[}CAL{]}}}\label{cm.7-why-copper-conducts-so-well-fcal}}

The Drude relation \(\sigma = ne^2\tau/m_e\) with
\(\sigma_{\mathrm{Cu}} = 5.96\times 10^7\ \mathrm{S/m}\) gives a
relaxation time \(\tau = 2.5\times 10^{-14}\ \mathrm{s}\) and a mean
free path \(\ell = v_F\tau = 39\ \mathrm{nm}\). Copper's \textbf{filled}
\(3d^{10}\) shell hands its \(4s^1\) electron to a low-scattering sea
--- exactly the property that, in iron, is \emph{absent}. This is the
hinge of the whole chapter: \textbf{a filled inner shell yields
conduction; an unfilled one yields magnetism.}

\hypertarget{cm.8-the-magnetic-field-of-a-current-and-joule-loss-f}{%
\subsection{\texorpdfstring{CM.8 The magnetic field of a current --- and
Joule loss
\textbf{{[}F{]}}}{CM.8 The magnetic field of a current --- and Joule loss {[}F{]}}}\label{cm.8-the-magnetic-field-of-a-current-and-joule-loss-f}}

A drifting charge tilts its rotation axis by \(\sim v/c\) (§EM.4),
shedding a small transverse twist perpendicular to the drift;
propagating sideways it wraps the wire azimuthally as the magnetic field
\(B = \mu_0 I/(2\pi r)\), with circulation sense set by the rotation
handedness (the right-hand rule). The transverse twist has two fates:
the \textbf{returned} part is the reactive \(B\) field itself (present
even at zero loss, as in a superconductor), and the \textbf{un-returned}
fraction is dissipation (Joule loss). Thus \(B\) and resistive loss
share one origin --- the transverse twist of moving charge --- as its
returned versus un-returned parts; \(B\) itself is not the loss.

\hypertarget{cm.9-irons-magnetism-the-principle-fcalo}{%
\subsection{\texorpdfstring{CM.9 Iron's magnetism --- the principle
\textbf{{[}F{]}/{[}CAL{]}/{[}O{]}}}{CM.9 Iron's magnetism --- the principle {[}F{]}/{[}CAL{]}/{[}O{]}}}\label{cm.9-irons-magnetism-the-principle-fcalo}}

Magnetism is aligned rotation. A magnetic moment is the rotation of an
unpaired electron --- the same rotation that \emph{is} electric charge
(§EM.1) and helicity (§EM.5). For a material to be a permanent
(ferro)magnet, \textbf{two conditions must both hold}, and the pair
derives exactly which elements magnetize.

\hypertarget{cm.10-criterion-1-unpaired-d-electrons-this-is-why-not-copper-f}{%
\subsection{\texorpdfstring{CM.10 Criterion 1 --- unpaired
\emph{d}-electrons (this is why \textbf{not copper})
\textbf{{[}F{]}}}{CM.10 Criterion 1 --- unpaired d-electrons (this is why not copper) {[}F{]}}}\label{cm.10-criterion-1-unpaired-d-electrons-this-is-why-not-copper-f}}

A localized moment requires an \textbf{unfilled} \(d\)-shell with
unpaired electrons (Hund's rule). For \(3d^{n_d}\) the unpaired count is
\(n_d\) if \(n_d\le 5\) else \(10-n_d\), and the spin-only moment is
\(\mu = \sqrt{n(n+2)}\,\mu_B\) (\texttt{vp\_iron\_magnetism.py},
condition 1):

\begin{longtable}[]{@{}
  >{\raggedright\arraybackslash}p{(\columnwidth - 8\tabcolsep) * \real{0.2000}}
  >{\raggedright\arraybackslash}p{(\columnwidth - 8\tabcolsep) * \real{0.2000}}
  >{\raggedright\arraybackslash}p{(\columnwidth - 8\tabcolsep) * \real{0.2000}}
  >{\raggedright\arraybackslash}p{(\columnwidth - 8\tabcolsep) * \real{0.2000}}
  >{\raggedright\arraybackslash}p{(\columnwidth - 8\tabcolsep) * \real{0.2000}}@{}}
\toprule\noalign{}
\begin{minipage}[b]{\linewidth}\raggedright
element
\end{minipage} & \begin{minipage}[b]{\linewidth}\raggedright
config
\end{minipage} & \begin{minipage}[b]{\linewidth}\raggedright
unpaired \(d\)
\end{minipage} & \begin{minipage}[b]{\linewidth}\raggedright
\(\mu_{\text{so}}\,[\mu_B]\)
\end{minipage} & \begin{minipage}[b]{\linewidth}\raggedright
moment?
\end{minipage} \\
\midrule\noalign{}
\endhead
\bottomrule\noalign{}
\endlastfoot
Mn & \(3d^54s^2\) & 5 & 5.92 & yes \\
Fe & \(3d^64s^2\) & 4 & 4.90 & yes \\
Co & \(3d^74s^2\) & 3 & 3.87 & yes \\
Ni & \(3d^84s^2\) & 2 & 2.83 & yes \\
\textbf{Cu} & \(3d^{10}4s^1\) & \textbf{0} & \textbf{0.00} &
\textbf{no} \\
\end{longtable}

Copper's \(3d^{10}\) shell is full: zero unpaired \(d\)-electrons, no
localized moment. Its \(4s^1\) is a free conduction electron, not a
localized moment. \textbf{Copper fails condition 1 and cannot be
magnetic --- the same filled shell that makes it a superb conductor.}

\hypertarget{cm.11-criterion-2-positive-exchange-this-is-why-not-manganese-fcal}{%
\subsection{\texorpdfstring{CM.11 Criterion 2 --- positive exchange
(this is why \textbf{not manganese})
\textbf{{[}F{]}/{[}CAL{]}}}{CM.11 Criterion 2 --- positive exchange (this is why not manganese) {[}F{]}/{[}CAL{]}}}\label{cm.11-criterion-2-positive-exchange-this-is-why-not-manganese-fcal}}

Having moments is not enough; neighbouring \(d\)-rotations must align
\textbf{parallel} (ferromagnetic), not antiparallel (antiferromagnetic).
The Bethe--Slater criterion uses the ratio \(D/d\) = interatomic
distance / \(3d\)-orbital diameter: above a critical value
(\(\approx 1.5\)) the exchange integral is positive (parallel,
ferromagnetic); below it the exchange is negative (antiparallel,
antiferromagnetic) (\texttt{vp\_iron\_magnetism.py}, condition 2):

\begin{longtable}[]{@{}llll@{}}
\toprule\noalign{}
element & \(D/d\) & exchange & alignment \\
\midrule\noalign{}
\endhead
\bottomrule\noalign{}
\endlastfoot
Cr & 1.18 & \(J<0\) & antiparallel (antiferro) \\
Mn & 1.47 & \(J<0\) & antiparallel (antiferro) \\
Fe & 1.63 & \(J>0\) & parallel (ferro) \\
Co & 1.82 & \(J>0\) & parallel (ferro) \\
Ni & 1.98 & \(J>0\) & parallel (ferro) \\
\end{longtable}

Manganese has \textbf{five} unpaired electrons --- more than iron ---
yet \(D/d = 1.47 < 1.5\) gives \(J<0\), so its moments align
antiparallel and cancel: \textbf{manganese fails condition 2 and is
antiferromagnetic.} Cobalt (\(D/d = 1.82\)) sits near the exchange
maximum.

\hypertarget{cm.12-the-derivation-77-classification-f}{%
\subsection{\texorpdfstring{CM.12 The derivation --- 7/7 classification
\textbf{{[}F{]}}}{CM.12 The derivation --- 7/7 classification {[}F{]}}}\label{cm.12-the-derivation-77-classification-f}}

Applying both conditions reproduces the observed magnetic class of the
late \(3d\) series with no fitted magnetic parameter
(\texttt{vp\_iron\_magnetism.py}, summary):

\begin{longtable}[]{@{}llllll@{}}
\toprule\noalign{}
element & unpaired \(d\) & \(D/d\) & derived & observed & ✓ \\
\midrule\noalign{}
\endhead
\bottomrule\noalign{}
\endlastfoot
Cr & 5 & 1.18 & antiferro & antiferro & ✓ \\
Mn & 5 & 1.47 & antiferro & antiferro & ✓ \\
Fe & 4 & 1.63 & ferro & ferro & ✓ \\
Co & 3 & 1.82 & ferro & ferro & ✓ \\
Ni & 2 & 1.98 & ferro & ferro & ✓ \\
Cu & 0 & --- & non-magnetic & non-magnetic & ✓ \\
Zn & 0 & --- & non-magnetic & non-magnetic & ✓ \\
\end{longtable}

\textbf{7/7.} Copper is non-magnetic by condition 1, manganese is
antiferromagnetic by condition 2, and only iron, cobalt, and nickel
satisfy both --- the ``why Fe and not Cu'' the physics whitepaper left
as input is here derived.

\hypertarget{cm.13-curie-temperature-domains-and-nablacdotmathbf-b-0-fcal}{%
\subsection{\texorpdfstring{CM.13 Curie temperature, domains, and
\(\nabla\!\cdot\mathbf B = 0\)
\textbf{{[}F{]}/{[}CAL{]}}}{CM.13 Curie temperature, domains, and \textbackslash nabla\textbackslash!\textbackslash cdot\textbackslash mathbf B = 0 {[}F{]}/{[}CAL{]}}}\label{cm.13-curie-temperature-domains-and-nablacdotmathbf-b-0-fcal}}

Thermal rotation destroys alignment above the Curie temperature, where
\(k_BT_C\) is of order the exchange energy: Fe 1043 K, Co 1388 K
(highest, matching the strongest exchange), Ni 627 K
(\texttt{vp\_iron\_magnetism.py}). Below \(T_C\), exchange aligns
moments into \textbf{domains}; an external field moves the domain walls.
Because \(\mathbf B\) \emph{is} a rotation axis (§EM.4), it is
divergence-free, \(\nabla\!\cdot\mathbf B = 0\): the poles are the two
ends of one axis and cannot be separated --- no monopoles, and cutting a
magnet yields two fresh N/S pairs. Hysteresis is the part of the
alignment that does not return --- the permanent magnet --- and heating
above \(T_C\) erases it. The \(\nabla\!\cdot\mathbf B = 0\) seen
abstractly in Chapter EM becomes, in iron, the macroscopic fact of
domains and permanent magnets. (The metallic moments are non-integer ---
Fe 2.22, Co 1.72, Ni 0.62 \(\mu_B\) --- an itinerant/band effect marked
\textbf{{[}O{]}}; the localized Hund count gives the correct
\emph{ordering} and the magnetic \emph{class}.)

\hypertarget{cm.14-the-unifying-picture}{%
\subsection{CM.14 The unifying
picture}\label{cm.14-the-unifying-picture}}

Copper and iron are the two fates of the \(d\)-shell:

\begin{itemize}
\tightlist
\item
  \textbf{Copper --- filled \(3d^{10}\) → free \(4s\) electron →
  conduction.} Electricity passes through copper as the longitudinal
  (\(\chi\to 0\)) electromagnetic mode of that free sea, while light
  (transverse) is reflected by it.
\item
  \textbf{Iron --- unfilled \(3d^6\) → unpaired moments → magnetism.}
  With positive exchange (\(D/d>1.5\)) the moments align into a
  ferromagnet; the aligned rotation is the magnetic field.
\end{itemize}

A filled inner shell makes a conductor; an unfilled one makes a magnet.
Both are the same rotation --- the rotation that is charge, light, and
helicity --- expressed as a delocalized current in copper and as an
aligned moment in iron.

\hypertarget{cm.15-falsification-and-reproducibility-f}{%
\subsection{\texorpdfstring{CM.15 Falsification and reproducibility
\textbf{{[}F{]}}}{CM.15 Falsification and reproducibility {[}F{]}}}\label{cm.15-falsification-and-reproducibility-f}}

\begin{itemize}
\tightlist
\item
  If a metal's plasma frequency did not predict its reflectivity edge
  (shiny below \(\omega_p\), transparent above), the plasma-reflection
  account of conduction-vs-light fails.
\item
  If manganese (more unpaired electrons than iron) were ferromagnetic,
  or copper (\(3d^{10}\)) were ferromagnetic, the two-condition
  criterion fails.
\item
  If the derived magnetic class disagreed with any late-\(3d\) metal,
  the derivation fails (it is 7/7).
\end{itemize}

\begin{verbatim}
python3 code/foundation/vp_copper_conduction.py   # plasma freq, skin depth, drift/signal/Fermi
python3 code/foundation/vp_iron_magnetism.py      # two-condition ferromagnetic criterion (7/7)
\end{verbatim}

Both modules are deterministic (stdout sha256 identical on re-run) and
standard-library only.

\emph{Copper and iron close the loop opened in Chapter EM: charge is
rotation, and that one rotation appears as a longitudinal current
threading copper's filled-shell electron sea, and as an aligned moment
in iron's unfilled \(d\)-shell. The conductor and the magnet are the
same physics, divided by whether the inner shell is full.}

\hypertarget{the-color-of-molecules-from-d-orbitals-to-paint-pigments}{%
\section{\texorpdfstring{The Color of Molecules: From \emph{d}-Orbitals
to Paint
Pigments}{The Color of Molecules: From d-Orbitals to Paint Pigments}}\label{the-color-of-molecules-from-d-orbitals-to-paint-pigments}}

\textbf{VP Chemistry \& Electromagnetism Whitepaper --- Chapter CC} DOI
(reserved): 10.5281/zenodo.20680541 · Builds on Chapter EM (light),
Chapter CH §CH.5 (crystal field), and Chapter CM. Numbers reproduced by
\texttt{vp\_molecular\_color.py} (and \texttt{vp\_crystal\_field.py});
deterministic, standard-library, self-checking.

\begin{quote}
Color is a subtraction. White light arrives; the molecule removes the
photons whose energy matches an electronic gap, and what is left ---
transmitted or reflected --- is the color we see. Four different gaps
give the four families of molecular color, but the rule is one: a gap
the size of a visible photon makes a colored substance.
\end{quote}

\hypertarget{cc.0-the-question}{%
\subsection{CC.0 The question}\label{cc.0-the-question}}

Why is copper sulfate blue, a carrot orange, cadmium-yellow paint
yellow, and titanium-white paint white? This chapter derives the color
of molecules and pigments from the electromagnetism of Chapter EM:
visible light is the transverse electromagnetic wave (§EM.6), and
\textbf{color is which visible photons a substance absorbs}, set by the
size of an electronic energy gap.

\hypertarget{cc.1-the-principle-color-is-selective-absorption-f}{%
\subsection{\texorpdfstring{CC.1 The principle --- color is selective
absorption
\textbf{{[}F{]}}}{CC.1 The principle --- color is selective absorption {[}F{]}}}\label{cc.1-the-principle-color-is-selective-absorption-f}}

A substance has color when its electronic structure absorbs \emph{part}
of the visible spectrum (400--700 nm). An absorbed photon (a transverse
light quantum) lifts an electron across an electronic \textbf{gap}; the
gap size fixes which wavelength is absorbed, and the remaining light is
the perceived color. For dissolved molecules the eye sees the
\textbf{complementary} of the absorbed band; for opaque pigments it sees
the \textbf{reflected} band. Two limits frame the rest: a gap
\emph{larger} than any visible photon absorbs nothing (white or
colorless), and a gap \emph{smaller} than every visible photon absorbs
everything (black).

\hypertarget{cc.2-energy-and-wavelength-f}{%
\subsection{\texorpdfstring{CC.2 Energy and wavelength
\textbf{{[}F{]}}}{CC.2 Energy and wavelength {[}F{]}}}\label{cc.2-energy-and-wavelength-f}}

The bridge is the photon relation of Chapter EM,
\[ E = \frac{hc}{\lambda}, \qquad \lambda[\mathrm{nm}] = \frac{1240}{E[\mathrm{eV}]}. \]
Visible light spans \(E = 1.77\ \mathrm{eV}\) (red, 700 nm) to
\(3.10\ \mathrm{eV}\) (violet, 400 nm). A gap in this energy window
produces color; outside it, the substance is white/colorless (gap too
large) or black (gap too small).

\hypertarget{cc.3-mechanism-1-dd-transitions-transition-metal-complexes-fcalo}{%
\subsection{\texorpdfstring{CC.3 Mechanism 1 --- \emph{d}--\emph{d}
transitions (transition-metal complexes)
\textbf{{[}F{]}/{[}CAL{]}/{[}O{]}}}{CC.3 Mechanism 1 --- d--d transitions (transition-metal complexes) {[}F{]}/{[}CAL{]}/{[}O{]}}}\label{cc.3-mechanism-1-dd-transitions-transition-metal-complexes-fcalo}}

In a transition-metal complex the ligands' electromagnetic \(1/r^2\)
field splits the geometric \(d\)-orbitals by the crystal-field gap
\(\Delta\) (Chapter CH §CH.5). An electron absorbs a photon to cross
\(\Delta\):
\[ \lambda_{\text{abs}}[\mathrm{nm}] = \frac{10^7}{\Delta[\mathrm{cm}^{-1}]}. \]
For \([\mathrm{Ti}(\mathrm H_2\mathrm O)_6]^{3+}\) (\(d^1\), single
transition), \(\Delta = 20300\
\mathrm{cm^{-1}}\) gives \(\lambda_{\text{abs}} = 493\ \mathrm{nm}\)
(blue-green absorbed) → reddish violet, matching observation
(\texttt{vp\_molecular\_color.py}, §1). The mechanism --- color
\emph{is} the \(d\)--\(d\) gap --- is forced by VP geometry
\textbf{{[}F{]}}; the spectrochemical ordering of \(\Delta\) is
calibration \textbf{{[}CAL{]}}; and for multi-electron ions (\(d^9\)
Cu²⁺, \(d^8\) Ni²⁺) the single-\(\Delta\) point model is insufficient
(broad multi-band absorption tails into the red, giving the real
blue/green), an honestly \textbf{{[}O{]}} limitation. Weak-field ligands
give a small \(\Delta\) (red absorption); strong-field ligands a large
\(\Delta\) (blue absorption).

\hypertarget{cc.4-mechanism-2-conjugated-pi-systems-organic-dyes-ff}{%
\subsection{\texorpdfstring{CC.4 Mechanism 2 --- conjugated \(\pi\)
systems (organic dyes)
\textbf{{[}F{]}/{[}F?{]}}}{CC.4 Mechanism 2 --- conjugated \textbackslash pi systems (organic dyes) {[}F{]}/{[}F?{]}}}\label{cc.4-mechanism-2-conjugated-pi-systems-organic-dyes-ff}}

In a conjugated polyene the \(\pi\)-electrons behave as a free electron
gas in a one-dimensional box of length \(L\) (the conjugation length).
The levels are \(E_n = n^2 h^2/(8 m_e L^2)\); with \(N = 2k\) electrons
(\(k\) double bonds) filling levels up to \(k\), the HOMO→LUMO gap is
\[ \Delta E = \frac{h^2}{8 m_e L^2}\,(N+1). \] Because \(L\) grows with
conjugation, \textbf{the gap shrinks and the absorption red-shifts} ---
the derived trend (\texttt{vp\_molecular\_color.py}, §2): short polyenes
(ethene, butadiene, hexatriene) absorb in the ultraviolet and are
colorless, while extended conjugation (β-carotene, 11 double bonds)
absorbs in the visible and is colored (orange; lycopene, red). The
scaling \(\Delta E \propto (N+1)/L^2\) and the UV→visible crossover are
forced \textbf{{[}F{]}}; the absolute wavelengths from the free-electron
model overestimate for long chains (bond alternation caps the gap), so
the model's numbers are trend-only \textbf{{[}F?{]}}. This is why
carrots and tomatoes are colored: long \(\pi\)-conjugation pulls the gap
down into the visible.

\hypertarget{cc.5-mechanism-3-the-band-gap-inorganic-pigments-and-paint-fcal}{%
\subsection{\texorpdfstring{CC.5 Mechanism 3 --- the band gap (inorganic
pigments and paint)
\textbf{{[}F{]}/{[}CAL{]}}}{CC.5 Mechanism 3 --- the band gap (inorganic pigments and paint) {[}F{]}/{[}CAL{]}}}\label{cc.5-mechanism-3-the-band-gap-inorganic-pigments-and-paint-fcal}}

A solid pigment is a semiconductor with a band gap \(E_g\). Photons with
\(E > E_g\) (wavelength below the absorption edge
\(\lambda_{\text{edge}} = hc/E_g\)) are absorbed; photons with
\(E < E_g\) (longer wavelength) are reflected, and \textbf{the reflected
band is the color}. As \(E_g\) decreases, the edge sweeps across the
visible and the pigment runs \textbf{white → yellow → orange → red →
black} (\texttt{vp\_molecular\_color.py}, §3):

\begin{longtable}[]{@{}llll@{}}
\toprule\noalign{}
pigment & \(E_g\) (eV) & edge (nm) & reflected color \\
\midrule\noalign{}
\endhead
\bottomrule\noalign{}
\endlastfoot
TiO₂ titanium white & 3.05 & 407 & white (all visible reflected) \\
ZnO zinc white & 3.20 & 387 & white \\
As₂S₃ orpiment & 2.70 & 459 & yellow (violet/blue absorbed) \\
CdS cadmium yellow & 2.42 & 512 & yellow \\
CdS·Se cadmium orange & 2.10 & 590 & orange/red \\
HgS vermilion (cinnabar) & 2.00 & 620 & red \\
Fe₂O₃ red ochre (hematite) & 2.10 & 590 & red \\
CdSe cadmium red & 1.73 & 717 & red (only deep red reflected) \\
carbon black & \textasciitilde0.5 & --- & black (all visible
absorbed) \\
\end{longtable}

The same single rule \(\lambda_{\text{edge}} = hc/E_g\) orders the
entire cadmium pigment series and the historical mineral pigments
(vermilion, orpiment, hematite) \textbf{{[}F{]}}, with \(E_g\) a
measured material constant \textbf{{[}CAL{]}}. This is the physics of
paint: \textbf{the band gap is the color knob.} Titanium white reflects
everything because its gap (3.05 eV) sits just past the violet edge;
cadmium red absorbs nearly all visible because its gap (1.73 eV) sits at
the red edge.

\hypertarget{cc.6-mechanism-4-charge-transfer-intense-colors-fcal}{%
\subsection{\texorpdfstring{CC.6 Mechanism 4 --- charge transfer
(intense colors)
\textbf{{[}F{]}/{[}CAL{]}}}{CC.6 Mechanism 4 --- charge transfer (intense colors) {[}F{]}/{[}CAL{]}}}\label{cc.6-mechanism-4-charge-transfer-intense-colors-fcal}}

When a photon moves an electron between a metal and a ligand
(ligand-to-metal or metal-to-ligand charge transfer), the transition is
fully allowed and therefore \textbf{intense} --- 100--1000× stronger
than a \(d\)--\(d\) transition, so a trace gives a deep color.
Permanganate \(\mathrm{MnO_4^-}\) absorbs \(\sim 525\ \mathrm{nm}\)
(O→Mn charge transfer) and is deep purple; dichromate
\(\mathrm{Cr_2O_7^{2-}}\) is orange; Prussian blue is an Fe²⁺↔Fe³⁺
charge transfer (\texttt{vp\_molecular\_color.py}, §4). The mechanism is
the same electron-across-a-gap; only the gap's origin (metal↔ligand) and
its allowedness differ.

\hypertarget{cc.7-white-black-and-everything-between-f}{%
\subsection{\texorpdfstring{CC.7 White, black, and everything between
\textbf{{[}F{]}}}{CC.7 White, black, and everything between {[}F{]}}}\label{cc.7-white-black-and-everything-between-f}}

The two achromatic limits follow from the same rule and require no extra
assumption:

\begin{itemize}
\tightlist
\item
  \textbf{White}: gap larger than any visible photon
  (\(E_g > 3.1\ \mathrm{eV}\), edge in the ultraviolet). No visible
  light is absorbed; all is reflected. Titanium white, zinc white, table
  salt, snow.
\item
  \textbf{Black}: gap smaller than every visible photon
  (\(E_g < 1.8\ \mathrm{eV}\), edge in the infrared). All visible light
  is absorbed. Carbon black, graphite.
\item
  \textbf{Color}: gap inside the visible window. The substance absorbs
  the short-wavelength side of the edge and shows the long-wavelength
  side.
\end{itemize}

So white and black are not ``colors added'' but the two ends of where
the gap sits relative to the visible band.

\hypertarget{cc.8-the-unifying-picture}{%
\subsection{CC.8 The unifying picture}\label{cc.8-the-unifying-picture}}

All four families are one mechanism --- \textbf{a transverse visible
photon absorbed by lifting an electron across an electronic gap} ---
differing only in what sets the gap:

\begin{itemize}
\tightlist
\item
  \(d\)--\(d\): the crystal-field splitting \(\Delta\) (ligand
  electromagnetic \(1/r^2\) field).
\item
  conjugated \(\pi\): the HOMO--LUMO gap \(\propto (N+1)/L^2\)
  (conjugation length).
\item
  band gap: the semiconductor \(E_g\) (the pigment's electronic
  structure).
\item
  charge transfer: the metal↔ligand donor--acceptor gap.
\end{itemize}

The absorbing light is the transverse electromagnetic wave of Chapter EM
(\(\chi\to 90^\circ\)); only a photon whose energy matches the gap is
absorbed (resonance), and the gap size determines the wavelength and
hence the color. One principle --- gap-sized absorption --- and four
ways nature builds the gap.

\hypertarget{cc.9-falsification-and-reproducibility-f}{%
\subsection{\texorpdfstring{CC.9 Falsification and reproducibility
\textbf{{[}F{]}}}{CC.9 Falsification and reproducibility {[}F{]}}}\label{cc.9-falsification-and-reproducibility-f}}

\begin{itemize}
\tightlist
\item
  If a transition-metal complex's color did not track its crystal-field
  \(\Delta\) (weak-field red, strong-field blue absorption), the
  \(d\)--\(d\) account fails.
\item
  If extending \(\pi\)-conjugation did not red-shift absorption (UV for
  short, visible for long), the particle-in-a-box account fails.
\item
  If a semiconductor pigment's color did not follow its band-gap edge
  \(\lambda_{\text{edge}} = hc/E_g\) (large gap white, small gap black,
  intermediate colored), the band-gap account fails --- this is tested
  across the cadmium series and historical pigments.
\end{itemize}

\begin{verbatim}
python3 code/foundation/vp_molecular_color.py    # complementary color, d-d, FEM, band-gap pigments
python3 code/foundation/vp_crystal_field.py      # crystal-field Δ and 4/9
\end{verbatim}

Both modules are deterministic (stdout sha256 identical on re-run) and
standard-library only.

\emph{Color completes the optical arc of Chapter EM: light emerges from
the lattice, refraction and the rainbow split it by wavelength, and here
matter }subtracts* selected wavelengths by lifting electrons across gaps
--- crystal-field, conjugation, band-gap, and charge-transfer gaps ---
so that a carrot is orange and titanium-white is white for the same
reason, read through the size of a gap.*

\hypertarget{platinum-catalysis-and-the-electron-amplitude-wavelength-toward-water-purification}{%
\section{Platinum Catalysis and the Electron Amplitude Wavelength:
Toward Water
Purification}\label{platinum-catalysis-and-the-electron-amplitude-wavelength-toward-water-purification}}

\textbf{VP Chemistry \& Electromagnetism Whitepaper --- Chapter CT} DOI
(reserved): 10.5281/zenodo.20680541 · Builds on Chapter EM, Chapter CH
(§CH.5, §CH.11), and Chapter CM. Related to the VP Application
whitepaper (10.5281/zenodo.18043066: Pt catalyst / electrolysis /
desalination). Numbers reproduced by \texttt{vp\_platinum\_catalysis.py}
(deterministic, standard-library). Grades distinguish
\textbf{established} surface science ({[}CAL{]}/{[}F?{]}) from the
\textbf{VP-framework hypothesis} ({[}H{]}).

\begin{quote}
A catalyst is a matchmaker: it must hold a reagent firmly enough to
activate it, yet loosely enough to release the product. Platinum sits at
that balance for hydrogen. This chapter asks whether the balance can be
read as a \emph{resonance} of the electron's amplitude wavelength with
the metal surface --- and whether that picture helps purify water.
\end{quote}

\hypertarget{ct.0-the-question}{%
\subsection{CT.0 The question}\label{ct.0-the-question}}

Why is platinum such a good catalyst, and can the ``electron amplitude
wavelength'' of the VP framework illuminate the mechanism? The
motivation is practical: platinum catalyses the hydrogen-evolution
reaction in water electrolysis, which is central to producing clean
hydrogen and to purifying water --- the application this chapter
targets.

\hypertarget{ct.1-why-platinum-states-at-the-fermi-level-f}{%
\subsection{\texorpdfstring{CT.1 Why platinum --- states at the Fermi
level
\textbf{{[}F?{]}}}{CT.1 Why platinum --- states at the Fermi level {[}F?{]}}}\label{ct.1-why-platinum-states-at-the-fermi-level-f}}

Platinum is \([\mathrm{Xe}]4f^{14}5d^9 6s^1\): its \(5d\) shell has a
single hole, so there are \(d\)-states right at the Fermi level.
Electrons are therefore always available both to donate to an adsorbate
and to accept back-donation from it. This is the prerequisite for
surface catalysis, and it is exactly what gold (\(5d^{10}\), filled)
lacks --- a filled \(d\)-shell makes gold comparatively inert, the same
``filled versus unfilled \(d\)-shell'' hinge that separated conductor
from magnet in Chapter CM.

\hypertarget{ct.2-the-sabatier-principle-and-the-volcano-calf}{%
\subsection{\texorpdfstring{CT.2 The Sabatier principle and the volcano
\textbf{{[}CAL{]}/{[}F?{]}}}{CT.2 The Sabatier principle and the volcano {[}CAL{]}/{[}F?{]}}}\label{ct.2-the-sabatier-principle-and-the-volcano-calf}}

The Sabatier principle states that the best catalyst binds the key
intermediate at \emph{intermediate} strength: too weak and it cannot
activate the reagent, too strong and it cannot release the product (it
is poisoned). Plotting the hydrogen-evolution exchange current against
the hydrogen adsorption free energy \(\Delta G_{\mathrm H^*}\) gives a
\textbf{volcano} peaking at \(\Delta G_{\mathrm H^*}\approx 0\)
(\texttt{vp\_platinum\_catalysis.py}, §1):

\begin{longtable}[]{@{}
  >{\raggedright\arraybackslash}p{(\columnwidth - 6\tabcolsep) * \real{0.2500}}
  >{\raggedright\arraybackslash}p{(\columnwidth - 6\tabcolsep) * \real{0.2500}}
  >{\raggedright\arraybackslash}p{(\columnwidth - 6\tabcolsep) * \real{0.2500}}
  >{\raggedright\arraybackslash}p{(\columnwidth - 6\tabcolsep) * \real{0.2500}}@{}}
\toprule\noalign{}
\begin{minipage}[b]{\linewidth}\raggedright
metal
\end{minipage} & \begin{minipage}[b]{\linewidth}\raggedright
\(\Delta G_{\mathrm H^*}\) (eV)
\end{minipage} & \begin{minipage}[b]{\linewidth}\raggedright
\(\log_{10} j_0\) (A/cm²)
\end{minipage} & \begin{minipage}[b]{\linewidth}\raggedright
regime
\end{minipage} \\
\midrule\noalign{}
\endhead
\bottomrule\noalign{}
\endlastfoot
\textbf{Pt} & \textbf{−0.09} & \textbf{−3.1} & \textbf{peak
(optimal)} \\
Ir & −0.10 & −3.7 & good \\
Pd & −0.20 & −3.5 & good \\
Ni & −0.28 & −5.2 & good \\
Au & +0.30 & −5.4 & too weak \\
W & −0.55 & −5.9 & too strong \\
Hg & +0.50 & −12.6 & too weak \\
\end{longtable}

Platinum has \(\Delta G_{\mathrm H^*}\) closest to zero and the highest
exchange current --- it sits at the apex of the volcano. Tungsten and
molybdenum bind hydrogen too strongly (cannot release it); gold and
mercury too weakly (cannot activate it). The volcano data are
calibration \textbf{{[}CAL{]}}; the Sabatier balance is the established
organizing principle \textbf{{[}F?{]}}.

\hypertarget{ct.3-the-d-band-center-calf}{%
\subsection{\texorpdfstring{CT.3 The d-band center
\textbf{{[}CAL{]}/{[}F?{]}}}{CT.3 The d-band center {[}CAL{]}/{[}F?{]}}}\label{ct.3-the-d-band-center-calf}}

The Hammer--Nørskov d-band-center model refines the picture: adsorption
strength correlates with the mean energy \(\varepsilon_d\) of the
\(d\)-band relative to the Fermi level. A higher \(\varepsilon_d\)
(closer to the Fermi level) binds adsorbates more strongly. Platinum's
\(\varepsilon_d \approx -2.25\ \mathrm{eV}\) lands in the optimal
window, between the over-reactive early transition metals (high
\(\varepsilon_d\)) and the noble metals such as gold
(\(\varepsilon_d \approx -3.56\ \mathrm{eV}\), too low, weak binding)
(\texttt{vp\_platinum\_catalysis.py}, §2).

\hypertarget{ct.4-the-electron-amplitude-wavelength-a-vp-resonance-picture-h}{%
\subsection{\texorpdfstring{CT.4 The electron amplitude wavelength --- a
VP resonance picture
\textbf{{[}H{]}}}{CT.4 The electron amplitude wavelength --- a VP resonance picture {[}H{]}}}\label{ct.4-the-electron-amplitude-wavelength-a-vp-resonance-picture-h}}

The VP framework adds a complementary, geometric reading. The
catalytically active electron has a de Broglie \emph{amplitude
wavelength} \[ \lambda_e = \frac{h}{\sqrt{2 m_e E}}, \] where \(E\) is
the energy scale at which the electron engages the adsorbate ---
naturally the work function \(\varphi = 5.65\ \mathrm{eV}\) for
platinum. This gives \(\lambda_e = 5.16\ \mathrm{\AA}\), so the
\textbf{half-wavelength is \(\lambda_e/2 = 2.58\ \mathrm{\AA}\)},
commensurate with the platinum nearest-neighbour spacing
\(d_{\mathrm{Pt-Pt}} = 2.77\ \mathrm{\AA}\) --- a matching ratio of
\textbf{0.93} (\texttt{vp\_platinum\_catalysis.py}, §3).

The hypothesis: when the electron's half-wavelength fits the surface
geometry (one half-wave per surface bond), the active electron couples
\emph{resonantly} into the adsorbate's antibonding orbital, so
back-donation is enhanced and the adsorbate bond is weakened --- the
activation step of catalysis. Platinum's particular combination of work
function and lattice spacing places it inside this resonance window.
This is a VP-framework picture, graded \textbf{{[}H{]}}: it reproduces
the qualitative fact (platinum is exceptional) and supplies a geometric
\emph{why}, but it is a scale-matching argument, not a quantitative
activity prediction. \textbf{§CT.7 puts it to a direct quantum
simulation --- which does not support the simple geometric form, and
revises the picture accordingly.}

\hypertarget{ct.5-tuning-the-resonance-with-potential-the-electrolysis-lever-h}{%
\subsection{\texorpdfstring{CT.5 Tuning the resonance with potential ---
the electrolysis lever
\textbf{{[}H{]}}}{CT.5 Tuning the resonance with potential --- the electrolysis lever {[}H{]}}}\label{ct.5-tuning-the-resonance-with-potential-the-electrolysis-lever-h}}

The amplitude-wavelength picture suggests a control knob. An applied
electrode potential \(V\) shifts the effective energy of the active
electron by \(eV\), and therefore shifts \(\lambda_e\)
(\texttt{vp\_platinum\_catalysis.py}, §4): a cathodic bias raises the
electron energy and shortens the half-wavelength, an anodic bias
lengthens it. Because catalysis happens \emph{at} an electrode under
applied potential, the electrode potential is precisely the lever that
tunes the proposed resonance. At the thermodynamic water-splitting
potential of 1.23 V the half-wavelength moves to \(2.92\
\mathrm{\AA}\) --- a few percent off the lattice match --- which the
picture would read as a small detuning to be compensated by catalyst
geometry. This makes the hypothesis \emph{actionable} for electrode
design rather than merely descriptive.

\hypertarget{ct.6-application-water-splitting-and-purification-fcal}{%
\subsection{\texorpdfstring{CT.6 Application --- water splitting and
purification
\textbf{{[}F{]}/{[}CAL{]}}}{CT.6 Application --- water splitting and purification {[}F{]}/{[}CAL{]}}}\label{ct.6-application-water-splitting-and-purification-fcal}}

The target application is water. In electrolysis the platinum cathode
runs the hydrogen-evolution reaction
\(2\mathrm H_2\mathrm O + 2e^- \to \mathrm H_2 + 2\mathrm{OH}^-\) with
an overpotential of only \(\sim 30\ \mathrm{mV}\) --- essentially zero
--- versus \(\sim 300\ \mathrm{mV}\) for nickel
(\texttt{vp\_platinum\_catalysis.py}, §5), so the cathodic energy loss
is minimised. By Faraday's law a current of 1 A for one hour produces
\(18.66\ \mathrm{mmol}\) (\(\sim 418\ \mathrm{mL}\) at STP) of hydrogen
(\(n = 2\)).

For purification, electrolysis contributes three ways: (i) anodic
oxidation breaks down microbes and organic contaminants (disinfection);
(ii) the evolved hydrogen and reactive species reduce pollutants; and
(iii) combined with the application whitepaper's
electromagnetic/electric-field desalination, it forms an integrated
purification train. The thread back to this chapter: platinum's
catalytic efficiency --- read in the VP picture as
electron-amplitude-wavelength resonance --- is what makes the
low-voltage electrolysis (and hence the energy-efficient purification)
feasible.

\hypertarget{ct.7-simulation-test-of-the-resonance-hypothesis-vrefuted}{%
\subsection{\texorpdfstring{CT.7 Simulation test of the resonance
hypothesis
\textbf{{[}V{]}/refuted}}{CT.7 Simulation test of the resonance hypothesis {[}V{]}/refuted}}\label{ct.7-simulation-test-of-the-resonance-hypothesis-vrefuted}}

Rather than rest on the scale coincidence, the hypothesis is put to a
direct quantum simulation (\texttt{vp\_pt\_resonance\_sim.py}): a 1D
finite-difference Schrödinger model of a periodic platinum surface
(wells spaced \(d = 2.77\ \mathrm{\AA}\)) with an adsorbate site,
computing the adsorbate's local density of states --- the back-donation
capacity --- as a function of electron energy, hence wavelength, via the
retarded Green's function \(G = [(E+i\eta) - H]^{-1}\).

The result \textbf{does not support the simple geometric resonance}, and
the simulation says why:

\begin{enumerate}
\def\labelenumi{\arabic{enumi}.}
\tightlist
\item
  \textbf{\(\lambda_e = 2d\) is a Bragg \emph{band edge}, not a
  resonance.} A transfer-matrix band-structure calculation places the
  bottom of the first conduction band at \(\approx 4.5\ \mathrm{eV}\);
  the Bragg wavelength \(\lambda_e = 2d\) (\(E = 4.90\ \mathrm{eV}\))
  sits at that edge, with a gap below. The Bragg condition is where
  electrons are \emph{reflected}, not resonantly transmitted.
\item
  \textbf{The adsorbate LDOS does not peak at \(\lambda_e = 2d\)} ---
  there it is only \(\sim 2\%\) of its maximum. The simple
  ``half-wavelength fits the lattice spacing → enhanced back-donation''
  picture is therefore not borne out: geometric commensurability is not
  a dynamical resonance.
\item
  \textbf{What governs back-donation} is the metal density of states at
  the adsorbate energy together with the adsorbate's own resonance ---
  i.e.~the established Newns--Anderson / d-band-center picture, not a
  geometric wavelength match. Platinum is exceptional because its partly
  filled \(5d\) band puts states at the Fermi level \emph{and} its
  hydrogen binding sits at the Sabatier optimum, exactly the §CT.2--CT.3
  account.
\item
  \textbf{Where a wavelength effect legitimately enters is quantum
  confinement.} Sweeping the slab thickness \(M\) (the confinement
  length \(L = Md\)) at a fixed Fermi energy, the adsorbate LDOS
  \emph{oscillates} strongly (coefficient of variation \(\approx 0.84\),
  several reversals) --- the real quantum-well-state effect known to
  modulate the reactivity of thin metal films. The wavelength that
  matters is matched to the film \emph{thickness}, not to the in-plane
  Pt--Pt spacing.
\end{enumerate}

So the simulation refines the chapter rather than confirming it: the
\(\lambda_e/2 \approx d_{\mathrm{Pt-Pt}}\) coincidence of §CT.4 is a
coincidence, not a back-donation resonance.

\hypertarget{ct.8-is-reproduction-impossible-energy-resonance-reproduces-platinum-v}{%
\subsection{\texorpdfstring{CT.8 Is reproduction impossible? --- energy
resonance reproduces platinum
\textbf{{[}V{]}}}{CT.8 Is reproduction impossible? --- energy resonance reproduces platinum {[}V{]}}}\label{ct.8-is-reproduction-impossible-energy-resonance-reproduces-platinum-v}}

The right question, having refuted the geometric form, is \emph{whether
platinum's catalysis can be reproduced at all, or whether the failure is
fundamental.} It is not fundamental: the failure was a \textbf{wrong
choice of resonance variable}, and the correct one reproduces platinum
(\texttt{vp\_pt\_reproduction.py}).

The catalytic resonance is in \textbf{energy}, not geometry: the
adsorbate orbital \(\varepsilon_a\) resonates with the metal's
\textbf{d-band centre} \(\varepsilon_d\) (the Newns--Anderson picture).
A Newns--Anderson simulation with a semi-elliptical d-band self-energy
\(\Sigma = V^2 g_d\) shows the sharp adsorbate level broadening into a
resonance on the d-band --- the energy-resonance mechanism. Crucially,
the \textbf{d-band centre predicts the hydrogen binding}: across Ni, Rh,
Pd, Ir, Pt, Cu, Au, the (independently determined) \(\varepsilon_d\)
correlates with the measured \(\Delta G_{\mathrm H^*}\) at \(r = -0.91\)
(a non-circular test, since \(\varepsilon_d\) comes from band structure,
not from the binding). Reconstructing the Sabatier volcano (activity
\(\propto \exp(-|\Delta G_{\mathrm H^*}|/kT)\)) then places
\textbf{platinum (with iridium) at the apex} --- the experimental
result, reproduced from the electronic structure.

So the answer to ``why couldn't we reproduce it, and is it impossible?''
is:

\begin{itemize}
\tightlist
\item
  \textbf{The geometric resonance failed because it is genuinely wrong}
  --- \(\lambda_e = 2d\) is a Bragg band edge in any dimension and any
  model (robust solid-state physics), not a limitation of the 1D
  simulation.
\item
  \textbf{Platinum's catalysis is fully reproducible} --- but through
  the \emph{energy} resonance
  (\(\varepsilon_a \leftrightarrow \varepsilon_d\)), which reproduces
  the binding (\(r=-0.91\)) and the volcano apex (Pt).
\item
  \textbf{The ``electron amplitude wavelength'' must be reinterpreted}:
  not the de Broglie wavelength matched to the lattice constant, but the
  \textbf{d-electron energy} matched to the adsorbate orbital. Read that
  way, the resonance intuition is correct and platinum is reproduced.
\end{itemize}

The catalytic essence of platinum is therefore the \textbf{energy} of
the d-states that its one-hole-short \(5d\) shell places at the Fermi
level --- exactly the §CT.1--CT.3 account, now shown to be a resonance,
just an energetic one.

\hypertarget{ct.9-assessment-of-the-possibility-hreframed}{%
\subsection{\texorpdfstring{CT.9 Assessment of the possibility
\textbf{{[}H{]}→reframed}}{CT.9 Assessment of the possibility {[}H{]}→reframed}}\label{ct.9-assessment-of-the-possibility-hreframed}}

Honestly weighed across both simulations:

\begin{itemize}
\tightlist
\item
  \textbf{The resonance intuition was right; the variable was wrong.}
  Catalysis \emph{is} a resonance --- of the adsorbate orbital with the
  d-band (energy), not of the electron wavelength with the lattice
  (geometry). The energy version reproduces platinum at the volcano apex
  (\(r=-0.91\)); the geometric version is a band edge.
\item
  \textbf{What survives of the amplitude-wavelength idea.} Two things:
  (i) the correct reading of ``amplitude wavelength'' as d-electron
  \emph{energy}, which is the established and predictive d-band-centre
  descriptor; and (ii) a genuine, simulable \emph{quantum-confinement}
  effect (thin-film quantum-well states, §CT.7) by which electron
  wavelength modulates reactivity through film thickness --- an
  actionable lever for the water-splitting electrodes of §CT.6.
\item
  \textbf{Verdict.} Reproduction is not impossible; it is achieved here
  through the energy resonance. The honest, predictive descriptor is the
  d-band centre, and the constructive engineering levers are the d-band
  position (alloying, strain) and quantum confinement (film thickness,
  potential) --- both tied directly to designing efficient catalyst
  layers for water electrolysis and purification.
\end{itemize}

This is the value of pressing past ``refuted'' to ``why, and is it
possible'': the geometric coincidence is discarded, the energy resonance
is shown to reproduce platinum, and the amplitude-wavelength idea is
corrected into its working form rather than abandoned.

\hypertarget{ct.10-falsification-and-reproducibility-fv}{%
\subsection{\texorpdfstring{CT.10 Falsification and reproducibility
\textbf{{[}F{]}/{[}V{]}}}{CT.10 Falsification and reproducibility {[}F{]}/{[}V{]}}}\label{ct.10-falsification-and-reproducibility-fv}}

\begin{itemize}
\tightlist
\item
  Established: if the hydrogen-evolution volcano did not peak near
  \(\Delta G_{\mathrm H^*} = 0\), or if platinum were not near that
  peak, the Sabatier account fails.
\item
  Hypothesis (now tested): the §CT.4 amplitude-wavelength resonance
  predicted that the adsorbate back-donation capacity would peak when
  \(\lambda_e/2\) matches the Pt--Pt spacing. The simulation of §CT.7
  finds the opposite --- \(\lambda_e = 2d\) is a band edge and the
  adsorbate LDOS does not peak there --- so this specific form is
  \textbf{refuted}, and the surviving, testable claim is the quantum-
  confinement one: thin-film reactivity should oscillate with layer
  thickness (and with \(E_F\) under applied potential).
\end{itemize}

\begin{verbatim}
python3 code/foundation/vp_platinum_catalysis.py   # volcano, d-band, λ_e scale-matching, electrolysis
python3 code/foundation/vp_pt_resonance_sim.py     # geometric resonance test: band edge (refuted)
python3 code/foundation/vp_pt_reproduction.py      # energy resonance: reproduces Pt at volcano apex
\end{verbatim}

All three modules are deterministic (stdout sha256 identical on re-run)
and standard-library only.

\emph{Platinum extends the metals thread of Chapter CM --- filled \(d\)
gives conduction (copper), unfilled \(d\) gives magnetism (iron), and a
\(d\)-shell one hole short of full gives catalysis (platinum). The VP
reading proposed a resonance of the electron's amplitude wavelength with
the surface. Put to a quantum simulation, the }geometric* form fails
(the Bragg match is a band edge), but pressing on to ask whether
reproduction is possible shows that it is: the \emph{energy} resonance
--- the adsorbate orbital with the d-band centre --- reproduces platinum
at the volcano apex (\(r=-0.91\)). The resonance intuition was right;
the variable was energy, not geometry. What this buys for water
electrolysis and purification is concrete: the levers are the d-band
position and quantum confinement of thin catalyst layers, both tuning
the Fermi-level d-states whose energy is the true seat of platinum's
catalysis.*

\hypertarget{applications-catalysis-fertilizer-water-energy-and-materials}{%
\section{Applications: Catalysis, Fertilizer, Water, Energy, and
Materials}\label{applications-catalysis-fertilizer-water-energy-and-materials}}

\textbf{VP Chemistry \& Electromagnetism Whitepaper --- Chapter CA} DOI
(reserved): 10.5281/zenodo.20680541 · Companion to Chapters CH
(chemistry), CM (metals), CT (catalysis), EM (electromagnetism). Status
grades: \textbf{{[}F{]}} forced (zero free parameters) ·
\textbf{{[}F?{]}} forced-candidate · \textbf{{[}CAL{]}} calibration
input (measured constant) · \textbf{{[}V{]}} simulation-measured ·
\textbf{{[}H{]}} hypothesis · \textbf{{[}O{]}} open / refuted. Every
numeric claim is reproduced by a deterministic, standard-library module
(2× sha256 identical); module and case-ledger anchors are given inline.

\begin{quote}
This chapter carries the metals account (CM/CT) into five practical
domains: \textbf{catalysts, fertilizer, water purification, energy
storage, and new materials}. Its method is the one that already
succeeded for platinum (Chapter CT): where a speculative
\emph{amplitude/geometry} mechanism was proposed, it is tested; what
does not survive is reported as refuted and \textbf{re-grounded on
validated physics}; what does survive is kept and labelled. The result
is a single thread --- the \textbf{d-band energy resonance} --- running
through catalysis, fertilizer, and water electrolysis, with honest
negative verdicts for two device claims that do not hold.
\end{quote}

\hypertarget{ca.0-the-central-method-re-grounding-not-decoration}{%
\subsection{CA.0 The central method: re-grounding, not
decoration}\label{ca.0-the-central-method-re-grounding-not-decoration}}

The prior application materials built catalysis, desalination, and an
energy device on an \emph{electron amplitude} variable (fm-scale).
Chapter CT showed by direct simulation that this variable is the wrong
one: the geometric/amplitude resonance is a Bragg band edge (refuted),
and platinum is reproduced only by the \textbf{energy} resonance --- the
d-band center (\(\varepsilon_d\)) coupling to the adsorbate level
(Newns--Anderson / Hammer--Nørskov). This chapter applies that verdict
everywhere: each application is rebuilt on the validated mechanism, and
any device claim whose mechanism fails testing is reported as
\textbf{refuted} or \textbf{{[}H{]}} with a working alternative
supplied. Honesty about the two negative results (magnetic desalination;
the ESS ``amplitude → electricity'' generator) is not a footnote --- it
is the chapter's spine.

\hypertarget{ca.1-the-d-band-engine-one-variable-for-three-reactions-vf}{%
\subsection{\texorpdfstring{CA.1 The d-band engine --- one variable for
three reactions
\textbf{{[}V{]}/{[}F?{]}}}{CA.1 The d-band engine --- one variable for three reactions {[}V{]}/{[}F?{]}}}\label{ca.1-the-d-band-engine-one-variable-for-three-reactions-vf}}

Module \texttt{vp\_dband\_catalysis.py} generalises the platinum result
into a reusable engine. A sharp adsorbate level broadens on coupling to
the metal d-band (Newns--Anderson) \textbf{{[}V{]}}; the d-band center
\(\varepsilon_d\) then \emph{predicts} binding. For hydrogen the
correlation with measured adsorption is
\[ r(\varepsilon_d,\ \Delta G_H) = -0.915 \quad\text{(strong, non-circular)} \;[\mathbf{F?}], \]
and the Sabatier volcano (activity peaked at \(\Delta G_H\approx 0\))
reproduces the \textbf{Pt apex}. The \emph{same engine} applied to
nitrogen adsorption gives a different optimum and the \textbf{Fe/Ru
apex} of the ammonia volcano. Crucially, the real tuning lever is the
\textbf{energy}, not the lattice spacing: shifting \(\varepsilon_d\) by
alloying or strain moves a metal along the volcano
(\texttt{vp\_dband\_catalysis.py} §4). This \emph{replaces} the refuted
geometric resonance with a design rule that is actually used for
electrodes. \textbf{Falsifiable:} HER apex at \(\Delta G_H\approx0\)
(Pt), ammonia apex at \(E_N\approx-0.5\) eV (Fe/Ru); break either
ordering and the d-band picture is discarded.

\hypertarget{ca.2-fertilizer-ammonia-synthesis-ff}{%
\subsection{\texorpdfstring{CA.2 Fertilizer --- ammonia synthesis
\textbf{{[}F{]}/{[}F?{]}}}{CA.2 Fertilizer --- ammonia synthesis {[}F{]}/{[}F?{]}}}\label{ca.2-fertilizer-ammonia-synthesis-ff}}

Nitrogen fixation \(\mathrm{N_2 + 3H_2 \rightleftharpoons 2NH_3}\)
underpins half the human food supply. Module
\texttt{vp\_ammonia\_synthesis.py} closes it with the two VP pillars ---
equilibrium (CH.8) and the d-band catalyst engine (CA.1).

\textbf{Equilibrium forces the conditions {[}F{]}.} With
\(\Delta H^\circ=-91.8\) kJ/mol (exothermic) and
\(\Delta S^\circ=-198.1\) J/mol·K (gas moles \(4\to2\)),
\(\Delta G^\circ=\Delta H-T\Delta S\) gives the crossover
\(T^{*} = \Delta H/\Delta S = 463\) K, where \(K=1.000\). Below it the
reaction is favoured but \(\mathrm{N_2}\) is kinetically inert; above it
the equilibrium yield collapses. The module solves the stoichiometric
equilibrium and reproduces the \textbf{Le Chatelier compromise}
quantitatively --- e.g. equilibrium \(\mathrm{NH_3}\) rises with
pressure (\(\Delta n=-2\)) and falls with temperature:

\begin{longtable}[]{@{}lllll@{}}
\toprule\noalign{}
condition & 1 atm & 100 atm & 200 atm & 400 atm \\
\midrule\noalign{}
\endhead
\bottomrule\noalign{}
\endlastfoot
500 K & 10.8\% & 76.3\% & 82.6\% & 87.3\% \\
700 K & 0.6\% & 29.1\% & 40.7\% & 52.4\% \\
\end{longtable}

This is exactly why industry runs ≈450 °C / 150--300 atm / Fe catalyst
--- a forced three-way compromise, not a tuned choice.

\textbf{Catalysis selects the metal {[}F?{]}.} The rate-limiting step is
dissociation of \(\mathrm{N\equiv N}\) (945 kJ/mol, among the strongest
bonds; the bond-energy sum reproduces \(\Delta H\approx-93\) kJ/mol,
sign and magnitude correct). The catalyst volcano in the
nitrogen-binding descriptor peaks at \textbf{Fe/Ru/Os} --- the
industrial Haber--Bosch catalysts --- with Mo/W too strong (N trapped)
and Ni/Pt too weak (no dissociation). The volcano comes from the d-band
energy of CA.1, one thread with catalysis and water-splitting. Absolute
rates, promoter (\(\mathrm{K_2O}\)/\(\mathrm{Al_2O_3}\)) effects, and
surface microkinetics remain \textbf{{[}O{]}}.

\hypertarget{ca.3-water-purification-i-electrolysis-electrodes-ff}{%
\subsection{\texorpdfstring{CA.3 Water purification I --- electrolysis
electrodes
\textbf{{[}F{]}/{[}F?{]}}}{CA.3 Water purification I --- electrolysis electrodes {[}F{]}/{[}F?{]}}}\label{ca.3-water-purification-i-electrolysis-electrodes-ff}}

Module \texttt{vp\_water\_electrolysis.py} carries the d-band engine to
clean hydrogen and electrochemical water treatment. Water splitting
\(\mathrm{2H_2O\to2H_2+O_2}\) has a reversible minimum
\(E^\circ_{\text{cell}}=1.229\) V (\(\Delta G=237\) kJ/mol H\(_2\)) and
a thermoneutral voltage \(V_{tn}=\Delta H/(nF)=1.481\) V; the difference
\(T\Delta S=48.6\) kJ/mol is supplied as heat \textbf{{[}F{]}}. The real
cell voltage decomposes as
\(V_{\text{cell}}=1.23+\eta_{\text{HER}}+\eta_{\text{OER}}+iR\), and the
single largest loss is the \textbf{anode (OER)}. The HER volcano (CA.1)
places the cathode near solved (Pt apex). The OER is bounded by a
\emph{fundamental} limit: the universal scaling relation
\(\Delta G_{*\mathrm{OOH}}-\Delta G_{*\mathrm{OH}}\approx 3.2\) eV ties
the intermediates, so the two middle steps cannot both be 1.23 eV and
the best achievable overpotential is
\[ \eta_{\text{OER}}^{\min} = 1.6 - 1.23 = 0.37\ \text{V} \quad[\mathbf{F?}], \]
with the apex at \textbf{RuO₂/IrO₂} (the industrial OER catalysts). This
0.37 V is \emph{why water-splitting is permanently expensive} on
single-site catalysts --- beating it requires breaking the scaling
relation (bifunctional / 3D sites), an honest \textbf{{[}O{]}}.
Faraday's law gives 18.66 mmol H\(_2\) per A·h \textbf{{[}F{]}} and a
thermoneutral voltage efficiency \(V_{tn}/V_{\text{cell}}\approx 81\%\).
Electrocoagulation, electrodialysis, and capacitive deionization (CDI)
extend the same electrode science to treatment and desalination --- all
charge/membrane/potential based, \textbf{{[}F?{]}/{[}CAL{]}}.

\hypertarget{ca.4-water-purification-ii-magnetic-desalination-an-honest-verdict-f-orefuted}{%
\subsection{\texorpdfstring{CA.4 Water purification II --- magnetic
desalination: an honest verdict \textbf{{[}F{]} ·
{[}O{]}/refuted}}{CA.4 Water purification II --- magnetic desalination: an honest verdict {[}F{]} · {[}O{]}/refuted}}\label{ca.4-water-purification-ii-magnetic-desalination-an-honest-verdict-f-orefuted}}

The prior materials proposed separating \(\mathrm{H_2O}\) from
\(\mathrm{Na^+/Cl^-}\) at 0.45 T by an \emph{amplitude} difference.
Module \texttt{vp\_magnet\_desalination.py} replaces the prose with the
actual magnetic force. Water is weakly diamagnetic
(\(\chi_v\approx-9.05\times10^{-6}\)); its magnetic energy in a 0.45 T
field is
\[ \frac{U_{\text{mag}}}{k_BT} \approx 5.3\times10^{-9} \quad[\mathbf{F}], \]
nine orders of magnitude below thermal energy --- thermal motion
overwhelms any magnetic ordering, so \textbf{no separation occurs}. To
reach \(U_{\text{mag}}\approx k_BT\) would need \(B\approx 6173\) T,
physically impossible for a purifier (cf.~\textasciitilde45 T
continuous, \textasciitilde1200 T pulsed-destructive). The Lorentz force
on flowing ions is real but \(\sim10^{-2}\,k_BT\) across a channel: a
field can stir (MHD) but cannot \emph{desalinate by amplitude}.
\textbf{Verdict: the amplitude-separation device is refuted {[}O{]}.}
The working alternatives are supplied with real numbers: reverse osmosis
(seawater osmotic pressure \(\pi=i\,cRT\approx 30\) bar, matching the
measured \textasciitilde27 bar; thermodynamic minimum \(\approx0.83\)
kWh/m³), electrodialysis, and CDI. This is the platinum pattern again
--- refute the wrong mechanism, keep the validated one.

\hypertarget{ca.5-energy-storage-black-copper-solar-ess-split-honestly-ff-halt}{%
\subsection{\texorpdfstring{CA.5 Energy storage --- black-copper solar
ESS, split honestly \textbf{{[}F{]}/{[}F?{]} ·
{[}H{]}→alt}}{CA.5 Energy storage --- black-copper solar ESS, split honestly {[}F{]}/{[}F?{]} · {[}H{]}→alt}}\label{ca.5-energy-storage-black-copper-solar-ess-split-honestly-ff-halt}}

Module \texttt{vp\_ess\_thermal.py} separates the device into a
\textbf{real storage part} and a \textbf{hypothetical generation part}.

\textbf{Storage is real and quantified.} A black-copper absorber
(\(\alpha\ge0.98\)) at 800 W/m² collects \textasciitilde4.7 kWh/day per
m². High-heat-capacity media (alumina/sand) store it; the heat capacity
is the Dulong--Petit \(3R\) per atom of Chapter CH.7 (\(c_p\approx1223\)
J/kg·K at the high-T limit; measured 880 J/kg·K at room temperature), so
a 1 m³ store holds \textbf{\textasciitilde46 kWh} of heat over a 118 K
rise \textbf{{[}F?{]}}. Vacuum insulation gives a thermal time constant
\(\tau=mc/(UA)\approx 5.4\) days \textbf{{[}F?{]}} --- ample for
overnight delivery.

\textbf{Generation is hypothesis, with a working alternative.} The
claimed ``rotation-amplitude → electric field'' generator has no
validated mechanism and is marked \textbf{{[}H{]}}. Heat-to-electricity
is bounded by Carnot, \(\eta=1-T_c/T_h=28.2\%\) at the storage
temperature; realistic conversion uses a thermoelectric (Seebeck,
\(ZT=1\) → \textasciitilde5.5\%, \textasciitilde2.5 kWh) or a Stirling
engine (\textasciitilde½ Carnot → \textasciitilde14\%,
\textasciitilde6.5 kWh) \textbf{{[}F{]}}. The honest split --- real
thermodynamic storage, unproven generator replaced by a heat engine ---
is the same discipline as CA.4.

\hypertarget{ca.6-new-materials-black-copper-and-d-band-design-ff}{%
\subsection{\texorpdfstring{CA.6 New materials --- black copper and
d-band design
\textbf{{[}F{]}/{[}F?{]}}}{CA.6 New materials --- black copper and d-band design {[}F{]}/{[}F?{]}}}\label{ca.6-new-materials-black-copper-and-d-band-design-ff}}

Module \texttt{vp\_new\_materials.py} explains the black-copper absorber
by \emph{optics}, not amplitude, and ties materials design to the
d-band. Polished copper reflects (\(R\approx0.97\), absorptance
\(\alpha=0.03\)) because of its free-electron plasma (Chapter CM);
nanostructured/oxidised copper raises \(\alpha\) to \textasciitilde0.96
(a \textasciitilde32× gain) by graded-index impedance matching,
light-trapping, and intrinsic CuO absorption \textbf{{[}F?{]}} --- the
same electromagnetic physics as the copper plasma/skin effect, not an
fm-scale amplitude. A \textbf{selective absorber} then needs \emph{low}
infrared emittance: the Stefan--Boltzmann re-radiation
\(P=\varepsilon\sigma A(T^4-T_{\text{amb}}^4)\) at the storage
temperature is +704 W net for \(\varepsilon=0.05\) but \textbf{−377 W
for a blackbody surface} (\(\varepsilon=0.90\)) --- the latter
\emph{cannot even heat up} \textbf{{[}F{]}}. Thus low IR emittance is
essential, a concrete falsifiable design rule. Beyond coatings, the
d-band of CA.1 is the materials lever (alloying/strain tunes catalytic,
electronic, and magnetic properties --- the CM/CT trilogy), and crystal
packing fractions remain pure geometry (FCC/HCP
\(\pi/(3\sqrt2)=0.7405\), \textbf{{[}F{]}}, CH.13).

\hypertarget{ca.7-master-ledger-applications}{%
\subsection{CA.7 Master ledger ---
applications}\label{ca.7-master-ledger-applications}}

\textbf{{[}F{]} forced (zero free parameters):} ammonia
\(T^{*}=\Delta H/\Delta S\), Le Chatelier directions, bond-energy
\(\Delta H\) sign; water-splitting \(1.23\)/\(1.48\) V and Faraday
\(18.66\) mmol/A·h; magnetic energy \(\ll k_BT\) and the infeasible
\(B\); seawater osmotic pressure and desalination minimum energy;
Dulong--Petit storage, thermal time constant, Carnot/Seebeck/Stirling
efficiencies; Stefan--Boltzmann selective-absorber balance; crystal
packing fractions.

\textbf{{[}F?{]} forced-candidate:} \(\varepsilon_d\leftrightarrow\)
binding correlations and volcano apexes (HER→Pt, NH₃→Fe/Ru,
OER→RuO₂/IrO₂); d-band strain/alloy tuning; OER scaling-relation limit
0.37 V; absorber enhancement and selectivity.

\textbf{{[}V{]} simulation-measured:} Newns--Anderson adsorbate-DOS
broadening; volcano reconstructions.

\textbf{{[}CAL{]} calibration inputs:} \(\varepsilon_d\),
\(\Delta G_H\), \(E_N\), OER descriptors, \(\Delta H/\Delta S\), bond
energies, \(E^\circ\), susceptibility, emittance, insolation.

\textbf{{[}H{]} hypothesis → alternative:} the ESS ``amplitude →
electricity'' generator (replaced by a heat engine).

\textbf{{[}O{]} open / refuted:} magnetic amplitude-desalination
(refuted); absolute reaction rates and microkinetics; promoter effects;
scaling-relation circumvention.

\hypertarget{ca.8-falsification-criteria-f}{%
\subsection{\texorpdfstring{CA.8 Falsification criteria
\textbf{{[}F{]}}}{CA.8 Falsification criteria {[}F{]}}}\label{ca.8-falsification-criteria-f}}

\begin{itemize}
\tightlist
\item
  If the HER apex is not Pt (\(\Delta G_H\approx0\)) or the ammonia apex
  is not Fe/Ru (\(E_N\approx-0.5\) eV), the d-band volcano account is
  discarded.
\item
  If exothermic ammonia equilibrium does not fall with temperature and
  rise with pressure, the equilibrium account is discarded.
\item
  If a single-site OER catalyst stably shows \(\eta<0.3\) V, the
  scaling-relation limit needs revision.
\item
  If a 0.45 T field measurably desalinates seawater, the magnetic-energy
  analysis is wrong (it is not: \(U_{\text{mag}}/k_BT\sim10^{-9}\)).
\item
  If a non-selective (high-\(\varepsilon\)) black absorber maintains the
  storage temperature under the stated flux, the radiative-balance rule
  is wrong.
\end{itemize}

\hypertarget{ca.9-reproducibility}{%
\subsection{CA.9 Reproducibility}\label{ca.9-reproducibility}}

Six deterministic, standard-library modules and three case ledgers (29
rows) reproduce this chapter; \texttt{verify\_applications.py} checks
execution, determinism (2× sha256), standard-library-only dependency,
and ledger integrity in one pass (\textbf{6/6 PASS, 29 rows}).

\begin{verbatim}
python3 verify_applications.py        # all modules + ledgers
python3 vp_dband_catalysis.py         # HER→Pt, NH3→Fe/Ru, d-band tuning lever
python3 vp_ammonia_synthesis.py       # Haber compromise, T*, catalyst volcano
python3 vp_water_electrolysis.py      # 1.23/1.48 V, OER 0.37 V limit, Faraday
python3 vp_magnet_desalination.py     # magnetic separation refuted; RO/ED/CDI
python3 vp_ess_thermal.py             # thermal storage (CH.7); generation [H] + heat engine
python3 vp_new_materials.py           # black-copper optics; selective absorber; packing
\end{verbatim}

\emph{The five applications close on one idea carried from the metals
trilogy: the \textbf{d-electron energy} (not a lattice geometry, not an
amplitude) governs catalysis, and the same energy thread reaches
fertilizer and water-splitting; thermodynamics governs equilibrium and
storage; optics governs the absorber. Where a device mechanism could not
be validated it is refuted or marked hypothesis and a working
alternative is supplied. What is forced is falsifiable, what is measured
is labelled, what is open is admitted --- and all of it is
reproducible.}

\newpage

\hypertarget{application-co_2-electrocatalytic-reduction-the-d-band-co-binding-descriptor-ff}{%
\section{\texorpdfstring{Application: CO\(_2\) Electrocatalytic
Reduction --- the \(d\)-band CO-binding descriptor
\textbf{{[}F{]}/{[}F?{]}}}{Application: CO\_2 Electrocatalytic Reduction --- the d-band CO-binding descriptor {[}F{]}/{[}F?{]}}}\label{application-co_2-electrocatalytic-reduction-the-d-band-co-binding-descriptor-ff}}

\emph{Numbers reproduced by \texttt{vp\_co2\_reduction.py}
(deterministic, standard-library, 2× sha256). This section re-grounds
the carbon-dioxide problem on the verified \(d\)-band energy descriptor,
replacing the refuted electron-amplitude account (in which CO\(_2\)
``dissociated'' when an amplitude exceeded \(\sim 300\) fm \(>\sqrt2\)).
The amplitude figure had no measured basis; the controlling variable is
the \(^*\)CO adsorption energy, exactly as for the hydrogen and nitrogen
volcanoes (Chapter CT, §CA).}

\hypertarget{the-mechanism-selectivity-is-set-by-co-binding-f}{%
\subsection{\texorpdfstring{The mechanism --- selectivity is set by
\(^*\)CO binding
\textbf{{[}F?{]}}}{The mechanism --- selectivity is set by \^{}*CO binding {[}F?{]}}}\label{the-mechanism-selectivity-is-set-by-co-binding-f}}

Electrochemical CO\(_2\) reduction (CO\(_2\)RR) proceeds through
adsorbed intermediates, and its \emph{product selectivity} is governed
by a single descriptor: the binding energy of adsorbed CO,
\(\Delta E_{\mathrm{CO}}\) --- the same \(d\)-band quantity that orders
hydrogen (\(\Delta G_H\), apex Pt) and nitrogen (\(E_N\), apex Fe/Ru).
The rate-limiting first electron transfer to \(^*\)CO\(_2^{-}\) /
\(^*\)COOH / \(^*\)OCHO is what a catalyst must stabilize; how well it
does so tracks \(\Delta E_{\mathrm{CO}}\). The module finds
\(\Delta E_{\mathrm{CO}}\) correlated with the Hammer-- Nørskov
\(d\)-band centre \(\varepsilon_d\) at \(r = -0.97\) --- confirming the
descriptor and reusing the same engine, not a new one.

\hypertarget{thermodynamics-equilibria-near-0-v-onset-far-below-f}{%
\subsection{\texorpdfstring{Thermodynamics --- equilibria near 0 V,
onset far below
\textbf{{[}F{]}}}{Thermodynamics --- equilibria near 0 V, onset far below {[}F{]}}}\label{thermodynamics-equilibria-near-0-v-onset-far-below-f}}

The multi-electron product potentials (vs RHE) cluster near zero, yet
the observed onset is much more negative (\(-0.8\) to \(-1.1\) V)
because the CO\(_2^{\bullet-}\) radical anion lies near \(-1.9\) V:

\begin{longtable}[]{@{}lll@{}}
\toprule\noalign{}
Product & \(n\,(e^-)\) & \(E^\circ\) (V vs RHE) \\
\midrule\noalign{}
\endhead
\bottomrule\noalign{}
\endlastfoot
CO & 2 & \(-0.10\) \\
HCOOH (formate) & 2 & \(-0.12\) \\
CH\(_3\)OH & 6 & \(+0.03\) \\
C\(_2\)H\(_4\) & 12 & \(+0.08\) \\
CH\(_4\) & 8 & \(+0.17\) \\
(competing) H\(_2\) & 2 & \(0.00\) \\
\end{longtable}

These equilibria are forced from formation free energies by
\(E^\circ = -\Delta G^\circ/(nF)\); the numeric gate independently
reproduces \(E^\circ(\mathrm{CO}_2\!\to\!\mathrm{CO}) = -0.104\) V and
\(E^\circ(\mathrm{CO}_2\!\to\!\mathrm{CH}_4) = +0.169\) V from
\(\Delta G_f\) data.

\hypertarget{coppers-uniqueness-the-falsifiable-headline-f}{%
\subsection{\texorpdfstring{Copper's uniqueness --- the falsifiable
headline
\textbf{{[}F?{]}}}{Copper's uniqueness --- the falsifiable headline {[}F?{]}}}\label{coppers-uniqueness-the-falsifiable-headline-f}}

Sorting metals by \(\Delta E_{\mathrm{CO}}\) reproduces Hori's
four-group classification:

\begin{itemize}
\tightlist
\item
  \textbf{Strong binding} (Pt, Pd, Ni, Rh, Co, Ru): the surface is
  CO-poisoned and evolves H\(_2\).
\item
  \textbf{Intermediate binding} (\(\Delta E_{\mathrm{CO}} \approx -0.5\)
  eV): CO stays long enough to be reduced further to
  \textbf{hydrocarbons and alcohols} --- this is \textbf{copper, and
  copper alone} among pure metals.
\item
  \textbf{Weak binding} (Au, Ag, Zn): CO desorbs and is the product
  (CO-selective).
\item
  \textbf{\(sp\)-metals} (Sn, Pb, Bi, In, Cd): bind \(^*\)OCHO through
  oxygen and give \textbf{formate}.
\end{itemize}

That copper is the only pure metal producing significant C\(_2\)+
hydrocarbons is the central, established experimental fact (Hori),
reproduced here as a consequence of intermediate \(^*\)CO binding ---
the CO\(_2\)RR analogue of ``Pt is the HER apex.''

\hypertarget{the-scaling-relation-ceiling-f}{%
\subsection{\texorpdfstring{The scaling-relation ceiling
\textbf{{[}F?{]}}}{The scaling-relation ceiling {[}F?{]}}}\label{the-scaling-relation-ceiling-f}}

Because \(^*\)COOH, \(^*\)CO, and \(^*\)CHO all bind through carbon,
their energies scale together with a nearly constant offset
\(\Delta E_{\mathrm{CHO}} \approx \Delta E_{\mathrm{CO}} + 0.74\) eV.
The \(^*\)CO \(\to\) \(^*\)CHO step is therefore potential-limiting,
giving a methane limiting potential on copper of \(U_L \approx -0.74\) V
(Peterson--Nørskov). This \(0.74\) eV floor cannot be removed by tuning
\(\varepsilon_d\) alone --- a linear-scaling constraint that is the
direct cousin of the OER \(0.37\) V overpotential floor in the water
chapter. It is the reason CO\(_2\)RR is intrinsically difficult, stated
as forced geometry of the scaling lines, not as a fitted loss.

\hypertarget{grade-and-falsification}{%
\subsection{Grade and falsification}\label{grade-and-falsification}}

\textbf{Forced {[}F{]}:} the equilibrium potentials and electron counts.
\textbf{Candidate-forced {[}F?{]}:} the
\(\varepsilon_d \leftrightarrow \Delta E_{\mathrm{CO}}\) descriptor,
copper's uniqueness, the volcano, and the scaling ceiling.
\textbf{Calibration {[}CAL{]}:} \(\Delta E_{\mathrm{CO}}\),
\(\varepsilon_d\), the equilibrium potentials, and the scaling offset.
\textbf{Open {[}O{]}:} absolute current densities and the full C--C
coupling microkinetics. \textbf{Falsifier:} copper is the apex for
hydrocarbons and CO\(_2\)RR selectivity follows the \(^*\)CO-binding
order; if a pure non-copper metal produces more hydrocarbons at the same
overpotential, or the binding order is broken, the \(d\)-band CO
descriptor is refuted.

\newpage

\hypertarget{appendix-o-open-items-register-what-is-honestly-unsolved}{%
\section{Appendix O: Open-Items Register (what is honestly
unsolved)}\label{appendix-o-open-items-register-what-is-honestly-unsolved}}

This appendix consolidates every {[}O{]} / {[}H{]} item carried in this
volume, with the specific physics that would be required to close it.
None is a defect to be fixed; each is open research. Closing any of them
would be \emph{new physics}, not a correction --- and forcing a closure
would reintroduce exactly the kind of overfitting (the retracted
\(\sqrt{2.5}\) rule, the refuted electron amplitude) that this program
has been careful to discard.

\hypertarget{o.1-chemistry-peripheral-quantitative-opens-independently-re-examined}{%
\subsection{O.1 Chemistry --- peripheral quantitative opens
(independently
re-examined)}\label{o.1-chemistry-peripheral-quantitative-opens-independently-re-examined}}

Each was re-examined by \texttt{vp\_open\_items\_reexam.py}; all four
correctly remain open because the closing variable lies \emph{outside}
VP \(\pi\)-geometry.

\begin{longtable}[]{@{}
  >{\raggedright\arraybackslash}p{(\columnwidth - 4\tabcolsep) * \real{0.3333}}
  >{\raggedright\arraybackslash}p{(\columnwidth - 4\tabcolsep) * \real{0.3333}}
  >{\raggedright\arraybackslash}p{(\columnwidth - 4\tabcolsep) * \real{0.3333}}@{}}
\toprule\noalign{}
\begin{minipage}[b]{\linewidth}\raggedright
Open item
\end{minipage} & \begin{minipage}[b]{\linewidth}\raggedright
Module
\end{minipage} & \begin{minipage}[b]{\linewidth}\raggedright
Closing physics (outside VP geometry)
\end{minipage} \\
\midrule\noalign{}
\endhead
\bottomrule\noalign{}
\endlastfoot
Absolute bond-dissociation energy & \texttt{vp\_bond\_energy.py} & the
Morse range \(\beta = a\,r_e\) (anharmonicity; lone-pair / polarity). A
direct test found \(\beta\) uncorrelated with bond order
(\(r = +0.18\)), so \(\sqrt2\) geometry alone cannot close it. \\
Defect knockdown (real vs ideal strength) & \texttt{vp\_strength.py} &
dislocation density \(\rho_{\mathrm{disl}}\) (processing-dependent;
Taylor \(\tau=\alpha G b\sqrt{\rho}\)) --- intrinsically a calibration
input. \\
Naïve shell correction & \texttt{vp\_shell\_correction.py} & the nuclear
spin--orbit shell model (Mayer--Jensen); a uniform proximity bonus
over-corrects because \(\Delta B\) changes sign across magic numbers. \\
Electron affinity / absolute IE & \texttt{vp\_dblock\_chemistry.py} &
electron screening constants (Clementi--Raimondi); Slater rules give
only the correlation (sign correct \(8/18\)). \\
\end{longtable}

\hypertarget{o.2-electromagnetism-foundational-opens-deeper-than-the-chemistry-set}{%
\subsection{O.2 Electromagnetism --- foundational opens (deeper than the
chemistry
set)}\label{o.2-electromagnetism-foundational-opens-deeper-than-the-chemistry-set}}

\begin{longtable}[]{@{}
  >{\raggedright\arraybackslash}p{(\columnwidth - 4\tabcolsep) * \real{0.3333}}
  >{\raggedright\arraybackslash}p{(\columnwidth - 4\tabcolsep) * \real{0.3333}}
  >{\raggedright\arraybackslash}p{(\columnwidth - 4\tabcolsep) * \real{0.3333}}@{}}
\toprule\noalign{}
\begin{minipage}[b]{\linewidth}\raggedright
Open item
\end{minipage} & \begin{minipage}[b]{\linewidth}\raggedright
Grade
\end{minipage} & \begin{minipage}[b]{\linewidth}\raggedright
Status
\end{minipage} \\
\midrule\noalign{}
\endhead
\bottomrule\noalign{}
\endlastfoot
\textbf{Electromagnetic coupling \(\alpha_{\mathrm{em}}\) (the
``\(1/137\)'')} & {[}CAL{]} & The \emph{strength} of electromagnetism is
a measured input; the framework derives the \emph{form} (\(1/r^2\),
\(c\), \(E/B\), polarization, the angle law) but not the magnitude.
Physics §14.5 states it is not derived. A topological / phase-winding
route has been scaffolded but is computationally heavy and is \textbf{on
hold}; there is no known low-cost shortcut. Closing it is new
physics. \\
\textbf{Long-range / unscreened EM (global \(U(1)\))} & {[}H{]} &
Masslessness and infinite range rest on the breaking of a \emph{global}
(not gauged) \(U(1)\), giving a Goldstone mode --- a stated departure
from gauge theory. The AQD dynamics supply causality (§EM.12) but not
the gauge structure; this remains a hypothesis. \\
\textbf{Absolute vector (magnetic / radiative) sector} & {[}O{]} & The
longitudinal (Coulomb / \(E\)) sector --- including the dynamic
equations, Poynting flow, and causality --- is closed in §EM.12. The
full vector \(\mathbf E/\mathbf B\) curl structure (Faraday's
\(\nabla\times\mathbf E\) and Ampère--Maxwell's
\(\nabla\times\mathbf B\) in vector form, with the magnetic field as an
independent transverse pseudovector and its absolute magnitude) is
\textbf{not} supplied by the scalar/longitudinal treatment. This is the
principal open structural item of the EM chapter. \\
\textbf{\(\gamma\gamma\) genesis dynamics} & {[}O{]} & How charge is
created from photon collisions (\(\gamma\gamma\to e^+e^-\)) at the level
of lattice dynamics is open. \\
\end{longtable}

\hypertarget{o.3-what-is-closed-for-contrast}{%
\subsection{\texorpdfstring{O.3 What is \emph{closed} (for
contrast)}{O.3 What is closed (for contrast)}}\label{o.3-what-is-closed-for-contrast}}

The longitudinal electrodynamics is complete and reproducible: light
emergence (\(c^2=B/\rho\)), the static Coulomb \(1/r^2\), the
\(|\mathbf B|=(v/c)|\mathbf E|\) ratio and \(\nabla\!\cdot\mathbf B=0\),
polarization as helicity, the conduction--radiation angle, blackbody
radiation, refraction, and --- as of this edition --- the
\textbf{dynamic field equations, Poynting's theorem, energy
conservation, uniqueness, and light-cone causality} (§EM.12). On the
chemistry side the forced body is complete: bonding geometry,
coordination color, the metal \(d\)-shell fates, thermodynamics,
equilibrium, kinetics, solution and electrochemistry, materials
geometry, and the nuclear skeleton; and the applications are re-grounded
on the verified \(d\)-band descriptor (catalysis, ammonia, water
electrolysis, energy storage, new materials, CO\(_2\) reduction), with
the non-working mechanisms (magnetic desalination,
amplitude-to-electricity ESS) honestly refuted and replaced by working
alternatives.

\newpage

\hypertarget{conclusion}{%
\section{Conclusion}\label{conclusion}}

This volume completes a single arc from one anchor. \textbf{Light
emerges} from the jammed lattice as the unique surviving longitudinal
wave at \(c^2 = B/\rho\), shown in closed form and by a deterministic
simulation (0.06\%), which also recovers linear dispersion, the collapse
of the transverse branch, and the framework's \(-19\) ppm mass-ratio
residual (§EM.2). On that foundation \textbf{charge is rotation}, the
Coulomb \(1/r^2\) is three-dimensional flux conservation, \(E\) and
\(B\) are the longitudinal and transverse parts of one twist,
polarization is helicity, and one wave spans conduction and radiation by
a single angle (§EM.3--EM.8). This edition closes the \textbf{dynamic
and causal structure} of that wave in longitudinal form --- the
first-order field equations, the Poynting energy flow with exact
conservation, uniqueness, and a strict light cone with quantum
microcausality (§EM.12, from the AQD axiomatic dynamics).

\textbf{Chemistry follows by sphere geometry and rotation.} Points
\emph{on} a sphere give molecular shape (\(\arccos(-1/3)\)); spherical
harmonics under a ligand field give the crystal-field ratio \(4/9\);
sphere \emph{packing} gives \(0.7405\); lattice points \emph{inside} a
sphere give the nuclear magic skeleton. Rotation threads the rest ---
the heat-capacity staircase, \(\Delta G = \Delta H - T\Delta S\), the
Boltzmann tail, and rotational-dent transfer (the Daniell cell at 1.10
V). The metal chapters split the \(d\)-shell into conduction (copper)
and magnetism (iron) and into color (\(d\)--\(d\) transitions); the
catalysis chapter records, in full, how the early electron-amplitude
resonance picture was \textbf{refuted} (a Bragg band edge, not a
resonance) and \textbf{reproduced correctly} by the \(d\)-band energy
descriptor.

\textbf{The applications close the arc on verified physics.} A single
\(d\)-band engine explains hydrogen (apex Pt), nitrogen for ammonia
(apex Fe/Ru), and \textbf{CO\(_2\) reduction (the uniqueness of copper
for hydrocarbons)}; water electrolysis is bounded by the OER scaling
floor; energy storage is real in its Dulong--Petit heat capacity and
Carnot-bounded in its conversion; and the mechanisms that do \emph{not}
work --- magnetic desalination, amplitude-to-electricity generation ---
are refuted honestly and replaced by working alternatives. The same
discipline that discards a beloved but unworkable idea is what lets the
rest be built upon.

What is \textbf{forced} is falsifiable and stated with its criterion;
what is \textbf{measured} is labelled and never tuned per result; what
is \textbf{open} --- most fundamentally the electromagnetic coupling
\(\alpha_{\mathrm{em}}\), and the absolute vector sector of
electromagnetism --- is admitted in the Open-Items Register (Appendix O)
rather than disguised. The result is intended to be both substantially
more detailed than the original working notes and fully reproducible: a
38-module chemistry/EM core with a 176-entry case ledger and a 9-module
/ 46-entry applications package, verification harnesses (38/38 and 9/9),
and numeric-drift gates (both pass, zero drift) accompany this text, and
every number above can be regenerated by the reader.

\hypertarget{reproduction-summary}{%
\subsection{Reproduction (summary)}\label{reproduction-summary}}

\begin{verbatim}
# chemistry + EM core
python3 code/verify_chemistry.py --root .     # 38/38 modules pass; ledgers 114 + 62 pass
python3 code/gate_numeric_chem.py             # numeric drift gate: PASS (zero drift)
python3 code/foundation/vp_light_emergence.py # light emerges at c^2 = B/rho
python3 code/foundation/vp_em_field_dynamics.py  # dynamic field, Poynting, causality (EM.12)

# applications
python3 apps/verify_applications.py           # 9/9 modules pass; ledger 46 rows
python3 apps/vp_chem_numeric_ssot.py          # applications drift gate: PASS (17 anchors)
python3 apps/vp_co2_reduction.py              # CO2 reduction: copper uniqueness

sha256sum -c SHA256SUMS.txt                   # file integrity
\end{verbatim}

\hypertarget{data-and-code-availability}{%
\subsection{Data and code
availability}\label{data-and-code-availability}}

All code, data, and documents are released under open licenses
(documentation and data CC BY 4.0; code MIT) in the DOI package reserved
as 10.5281/zenodo.20680541, with checksums and a deterministic RunID.
The package is a supplement to the VP Theory whitepaper
(10.5281/zenodo.17932567), takes its field-dynamics addendum from the
AQD record (10.5281/zenodo.17423870), and supersedes the earlier
qualitative application notes (related: 10.5281/zenodo.18043066).

\end{document}
